\documentclass[12pt,a4paper]{report}

% ── Packages ──────────────────────────────────────────────────
\usepackage[T1]{fontenc}
\usepackage[utf8]{inputenc}
\usepackage[margin=2.5cm]{geometry}
\usepackage{lmodern}
\usepackage{microtype}
\usepackage{setspace}
\usepackage{titlesec}
\usepackage{titletoc}
\usepackage{tocloft}
\usepackage{fancyhdr}
\usepackage{xcolor}
\usepackage{graphicx}
\usepackage{float}
\usepackage{mdframed}
\usepackage{enumitem}
\usepackage{booktabs}
\usepackage{array}
\usepackage{tabularx}
\usepackage{multirow}
\usepackage{hyperref}
\usepackage{caption}
\usepackage{subcaption}
\usepackage{fontawesome5}
\usepackage{tcolorbox}
\usepackage{tikz}
\usepackage{amsmath}
\usepackage{parskip}

% ── Colors ────────────────────────────────────────────────────
\definecolor{brand}{HTML}{1A73E8}
\definecolor{brandlight}{HTML}{E8F0FE}
\definecolor{danger}{HTML}{EA4335}
\definecolor{success}{HTML}{34A853}
\definecolor{warning}{HTML}{FBBC04}
\definecolor{gray800}{HTML}{202124}
\definecolor{gray500}{HTML}{5F6368}
\definecolor{gray200}{HTML}{E8EAED}
\definecolor{dark}{HTML}{050510}
\definecolor{principleblue}{HTML}{0D47A1}
\definecolor{principlegreen}{HTML}{1B5E20}
\definecolor{principlered}{HTML}{B71C1C}
\definecolor{principlepurple}{HTML}{4A148C}
\definecolor{principleteal}{HTML}{004D40}
\definecolor{principleamber}{HTML}{E65100}

% ── tcolorbox styles ──────────────────────────────────────────
\tcbuselibrary{skins, breakable, raster}

\newtcolorbox{principleBox}[2][]{
  enhanced, breakable,
  colback=brandlight, colframe=brand,
  fonttitle=\bfseries\sffamily,
  title={#2},
  left=6pt, right=6pt, top=5pt, bottom=5pt,
  boxrule=1.2pt,
  arc=4pt,
  #1
}

\newtcolorbox{gmailIssueBox}[2][]{
  enhanced, breakable,
  colback=red!5!white, colframe=danger,
  fonttitle=\bfseries\sffamily,
  title={#2},
  left=6pt, right=6pt, top=5pt, bottom=5pt,
  boxrule=1.2pt,
  arc=4pt,
  #1
}

\newtcolorbox{lucidSolutionBox}[2][]{
  enhanced, breakable,
  colback=green!5!white, colframe=success,
  fonttitle=\bfseries\sffamily,
  title={#2},
  left=6pt, right=6pt, top=5pt, bottom=5pt,
  boxrule=1.2pt,
  arc=4pt,
  #1
}

\newtcolorbox{imageBox}[1]{
  enhanced,
  colback=gray200!40!white, colframe=gray200,
  fonttitle=\small\itshape,
  title={#1},
  left=4pt, right=4pt, top=4pt, bottom=4pt,
  boxrule=1pt, arc=4pt,
  halign=center, valign=center,
  height=6cm,
}

\newtcolorbox{keyBox}[1][]{
  enhanced,
  colback=gray800!8!white, colframe=gray800!30!white,
  left=6pt, right=6pt, top=4pt, bottom=4pt,
  boxrule=0.8pt, arc=3pt,
  #1
}

% ── Typography ────────────────────────────────────────────────
\usepackage{lmodern}
\renewcommand{\familydefault}{\sfdefault}

% ── Page style ────────────────────────────────────────────────
\setlength{\headheight}{15pt}
\pagestyle{fancy}
\fancyhf{}
\fancyhead[L]{\small\textcolor{gray500}{Lucid Mail — HCI Design Report}}
\fancyhead[R]{\small\textcolor{gray500}{\thepage}}
\fancyfoot[C]{\small\textcolor{gray200}{\rule{0.5\textwidth}{0.4pt}}}
\renewcommand{\headrulewidth}{0.4pt}
\renewcommand{\headrule}{\hbox to\headwidth{\color{gray200}\leaders\hrule height \headrulewidth\hfill}}

% ── Section styling ───────────────────────────────────────────
\titleformat{\chapter}[display]
  {\normalfont\huge\bfseries\color{gray800}}
  {\textcolor{brand}{\chaptertitlename\ \thechapter}}{16pt}{\Huge}
\titlespacing*{\chapter}{0pt}{0pt}{20pt}

\titleformat{\section}
  {\normalfont\Large\bfseries\color{gray800}\raggedright}
  {\textcolor{brand}{\thesection}}{0.8em}{}
\titleformat{\subsection}
  {\normalfont\large\bfseries\color{gray800}\raggedright}
  {\thesubsection}{0.6em}{}

% ── Custom commands ───────────────────────────────────────────
\newcommand{\principleTag}[1]{\texttt{\small\textcolor{brand}{[#1]}}}
\newcommand{\gmailBad}[1]{\textbf{\textcolor{danger}{#1}}}
\newcommand{\lucidGood}[1]{\textbf{\textcolor{success}{#1}}}
\newcommand{\featureRef}[1]{\textit{\textcolor{principlepurple}{#1}}}

% ── Hyperref ──────────────────────────────────────────────────
\hypersetup{
  colorlinks=true,
  linkcolor=brand,
  urlcolor=brand,
  citecolor=brand,
  pdftitle={Lucid Mail — HCI Design Principles Report},
  pdfauthor={HCI End Semester Project},
}

% ══════════════════════════════════════════════════════════════
\begin{document}
\onehalfspacing

% ── Title Page ───────────────────────────────────────────────
\begin{titlepage}
  \pagecolor{dark}
  \color{white}
  \centering
  \vspace*{3cm}

  {\fontsize{64}{72}\selectfont\bfseries
    \textcolor[HTML]{4285F4}{L}%
    \textcolor[HTML]{EA4335}{U}%
    \textcolor[HTML]{FBBC04}{C}%
    \textcolor[HTML]{34A853}{I}%
    \textcolor[HTML]{4285F4}{D}%
  }\\[0.4cm]

  {\LARGE\bfseries\color{white!90!dark} Mail}\\[1.2cm]

  {\large\color{white!60!dark}
    \rule{0.4\textwidth}{0.4pt}\\[0.5cm]
    HCI End-Semester Project\\
    \textbf{Design Principles Report}\\
    \rule{0.4\textwidth}{0.4pt}
  }\\[2.5cm]

  \begin{tcolorbox}[
    enhanced, sharp corners,
    colback=white!5!dark, colframe=white!15!dark,
    width=0.72\textwidth,
    left=18pt, right=18pt, top=14pt, bottom=14pt,
  ]
    \centering
    \color{white!80!dark}
    {\large\bfseries Abstract}\\[8pt]
    {\small
      Lucid Mail is a browser-based email client, built from the ground up as an
      Human-Computer Interaction study. Every design decision — from the resizable
      sidebar to the 30-second autosave interval — is anchored to a named,
      peer-reviewed HCI principle. This report traces each such decision through
      Shneiderman's 8 Golden Rules, Nielsen's 10 Heuristics, Norman's 6 Design
      Principles, Gestalt theory, Design Laws, Universal Design, web usability
      guidelines, and psychological persuasion frameworks.  A dedicated chapter
      then maps each design choice directly to a known usability issue in Gmail,
      demonstrating how principled design can close the gap between feature
      richness and cognitive ease of use.
    }
  \end{tcolorbox}

  \vfill
\end{titlepage}

  % ── Course & Team Details ──────────────────────────────────
  \clearpage
  \thispagestyle{empty}
\pagecolor{dark}
\color{white}
\vspace*{\fill}
\begin{center}
  \begin{tcolorbox}[
    enhanced, sharp corners,
    colback=white!5!dark, colframe=white!15!dark,
    width=0.65\textwidth,
    left=28pt, right=28pt, top=28pt, bottom=28pt,
  ]
    \centering\color{white!80!dark}
    {\large\bfseries\color{white} Course Details}\\[14pt]
    {\normalsize\color{white!85!dark} Human Computer Interaction (HCI)}\\[6pt]
    {\normalsize\color{white!55!dark} CS5015}\\[24pt]
    {\large\bfseries\color{white} Team Members}\\[14pt]
    {\normalsize\color{white!85!dark} Sudarshan S}\quad
    {\normalsize\color{white!45!dark} CS23B2007}\\[10pt]
    {\normalsize\color{white!85!dark} R Abinav}\quad
    {\normalsize\color{white!45!dark} ME23B1004}
  \end{tcolorbox}
\end{center}
\vspace*{\fill}
\clearpage
\pagecolor{white}
\color{black}

% ── ToC ───────────────────────────────────────────────────────
\tableofcontents
\clearpage

% ══════════════════════════════════════════════════════════════
\chapter{Introduction}
% ══════════════════════════════════════════════════════════════

\section{Project Overview}

Lucid Mail is a fully functional email web application built with React and TypeScript,
designed as a direct response to usability shortcomings in existing email clients —
most notably Gmail. Rather than iterating on an existing codebase, Lucid was designed
from first principles: every feature, layout choice, interaction pattern, and visual
decision was validated against a specific, named HCI rule before implementation.

The application provides:

\begin{itemize}[leftmargin=*, itemsep=2pt]
  \item A \textbf{Landing Page} with persuasion-driven onboarding, social proof, and a night-sky
        animation canvas to trigger visceral first impressions.
  \item A \textbf{Top Bar} containing the wordmark, a persistent search field with filters and
        recent-search memory, a zoom widget, settings, and a multi-account switcher.
  \item A \textbf{Resizable Sidebar} that snaps between compact icon mode and expanded text mode,
        housing navigation, custom labels, Meet links, and Chat contacts.
  \item An \textbf{Email List} with a contextual Welcome Header, skeleton loading states, rich
        empty states, and a context-sensitive bulk action bar.
  \item An \textbf{Email Detail} view with a two-column layout: an action sidebar and a main
        content area with smart replies, spam warnings, and context banners for Trash/Archive.
  \item A \textbf{Compose Window} with rich text editing, 30-second autosave, subject-empty
        warnings, recipient autocomplete, file attachments, and Undo Send.
  \item A full \textbf{Settings} page covering theme, density, interface zoom, notifications,
        customisable keyboard shortcuts, tags (custom labels), accounts, and email signature.
  \item A \textbf{Keyboard Shortcuts} system with a discoverable overlay and per-shortcut
        remapping in Settings.
  \item \textbf{Toast Notifications} with inline Undo actions.
  \item A \textbf{Confirm Dialog} for irreversible destructive operations.
\end{itemize}

\section{Design Philosophy}

The guiding philosophy of Lucid is captured in three words the landing page repeats: ``Email,
reimagined for how you think.'' Every decision subordinates visual novelty to cognitive ease.
Tesler's Law — that every system has irreducible complexity — drove us to absorb complexity on
the developer side rather than expose it to users. Zipf's Principle of Least Effort drove us
to make every common action (compose, archive, search) reachable in under two keystrokes.
And Raskin's Three Laws prevented us from ever designing a path that could silently lose a
user's data.

\begin{figure}[H]
  \centering
  \includegraphics[width=\linewidth]{mainpage.jpg}
  \caption{Lucid Mail application overview}
\end{figure}
% ══════════════════════════════════════════════════════════════
\chapter{Shneiderman's 8 Golden Rules}
% ══════════════════════════════════════════════════════════════

Ben Shneiderman's eight golden rules form the backbone of interaction design evaluation.
This chapter maps each rule to a specific, concrete feature of Lucid Mail.

% ─────────────────────────────────────────────────
\section{S1 — Strive for Consistency}

\begin{principleBox}{S1 · Strive for Consistency}
  Use identical sequences, layouts, terminology, and UI actions across every
  part of the system. Same error style, same button positions, same terminology.
\end{principleBox}

Consistency in Lucid operates at four levels:

\textbf{Visual vocabulary.} Every destructive action (Delete, Delete forever, Sign Out)
uses the same \textcolor{danger}{red} colour and the same icon family (Feather icons).
Every positive confirmation uses \textcolor{success}{green}. This rule is applied in the
Bulk Action Bar, the Email Detail action sidebar, the Accounts Settings tab, and the
Tags Settings tab — four separate components sharing one visual language.

\textbf{Button hierarchy.} Primary calls-to-action (Send, Create tag, Sign in) always use
the \texttt{btn-primary} class — blue background, white text, 7px/18px padding. Secondary
actions always use \texttt{btn-outline} — transparent background, border, muted text.
Destructive confirm buttons always inherit the danger colour. No exception was permitted.

\textbf{Terminology.} The application never calls the same thing by two names. The
navigation label ``Archive'' is identical to the toast message ``Archived'', the bulk bar
button text ``Archive'', and the keyboard shortcut description ``Archive selected.'' Users
who learn the term in one location recognise it in all others.

\textbf{Keyboard shortcut hints.} Everywhere a shortcut exists, it is shown in
parentheses in the tooltip — \texttt{Compose ( C )}, \texttt{Search mail ( / )},
\texttt{Keyboard shortcuts ( ? )}. The formatting pattern is always the same.


\begin{figure}[H]
  \centering
  \includegraphics[width=\linewidth]{inbox.jpg}
  \caption{Inbox Page}
\end{figure}

% ─────────────────────────────────────────────────
\section{S2 — Enable Frequent Users to Use Shortcuts}

\begin{principleBox}{S2 · Enable Shortcuts}
  Advanced users should bypass steps that slow them down;
  novices still follow the full path. Keyboard shortcuts, macros, autocomplete.
\end{principleBox}

Lucid implements a fully layered shortcut system, arguably the most thorough in any
student-built email client.

\textbf{Single-key shortcuts:} \texttt{C} (compose), \texttt{E} (archive selected),
\texttt{\#} (delete selected), \texttt{Z} (undo last action), \texttt{/} (focus search),
\texttt{?} (open shortcuts overlay).

\textbf{Two-key sequence shortcuts:} \texttt{G I} (go to Inbox), \texttt{G S} (go to Sent),
\texttt{G T} (go to Trash), \texttt{G D} (go to Drafts). These are debounced within 800 ms,
matching Gmail's own two-key navigation model (Jakob's Law, L8).

\textbf{Compose window shortcuts:} \texttt{Ctrl+B} (bold), \texttt{Ctrl+I} (italic),
\texttt{Ctrl+U} (underline), \texttt{Ctrl+Enter} (send). These mirror standard word-processor
conventions so existing muscle memory transfers immediately.

\textbf{Discovery without recall.} The \texttt{?} key opens a full overlay that lists every
shortcut, organised into three categories: Actions, Navigation, Search \& Help. The overlay
is also reachable from the Help icon in the Top Bar — satisfying both the expert (hotkey) and
the novice (visual discovery) paths simultaneously.

\textbf{Remappable shortcuts.} In Settings → Shortcuts, each shortcut has a ``badge'' that
can be clicked to capture a new key press. This is the only email client in this class that
implements per-shortcut remapping — directly supporting the Power Law of Practice (L5), because
users who have already built muscle memory on different clients can retrain Lucid to match.


\begin{figure}[H]
  \centering
  \includegraphics[width=\linewidth]{shortcut.jpg}
  \caption{Keyboard Shortcuts overlay (\texttt{?} key) with three categories}
\end{figure}

\begin{figure}[H]
  \centering
  \includegraphics[width=\linewidth]{set3.jpg}
  \caption{Settings → Shortcuts tab with remappable key badges}
\end{figure}


% ─────────────────────────────────────────────────
\section{S3 — Offer Informative Feedback}

\begin{principleBox}{S3 · Informative Feedback}
  Every action should have a visible, proportional system response.
  Minor actions: small feedback. Major actions: substantial feedback.
\end{principleBox}

Lucid applies proportional feedback across four tiers:

\textbf{Tier 1 — Micro feedback (hover/focus states).} Every button changes opacity or
background on hover. Every text input changes border colour to \texttt{brand-500} on focus
and returns to \texttt{border} on blur. This is implemented identically in the search bar,
compose recipient inputs, subject input, tag name input, and Settings text areas.

\textbf{Tier 2 — Status indicators.} The Compose Window title bar shows a live
\texttt{Saved} / \texttt{Saving…} indicator using a small \textcolor{success}{green}
checkmark, giving the user continuous awareness that their draft is safe — without interrupting
the writing flow. This responds to Raskin's First Law: never allow work to be lost.

\textbf{Tier 3 — Toast notifications.} Every significant action (archive, trash, send,
restore, permanently delete) produces a toast at the bottom of the screen with the action's
result in plain English: ``Archived'', ``Moved to Trash'', ``Message sent''. Toasts for
reversible actions include an ``Undo'' link. Toasts for irreversible ones (``Permanently
deleted'') carry no undo link — deliberately, because offering false hope of reversal
would violate user trust.

\textbf{Tier 4 — Loading skeletons.} When the email list is loading, eight skeleton rows
are rendered — animated grey bars matching the shape of real email rows. This uses the
Gestalt Closure principle (G1): the brain reads the familiar shape and understands ``data is
coming'' without any text instruction.



\begin{figure}[H]
  \centering
  \includegraphics[width=\linewidth]{red_saved.jpg}
  \caption{``Saved'' indicator in Compose window title bar}
\end{figure}


\begin{figure}[H]
  \centering
  \includegraphics[width=\linewidth]{feedback1.jpg}
  \caption{Toast notification with ``Trash'' action after archive}
\end{figure}


% ─────────────────────────────────────────────────
\section{S4 — Design Dialogs to Yield Closure}

\begin{principleBox}{S4 · Dialogs for Closure}
  Multi-step tasks should have a clear, recognisable end state.
  Users should never wonder ``did it work?''
\end{principleBox}

Closure is implemented in Lucid through \textit{context banners} — full-width inline strips
that appear at the top of the Email Detail view when the email's current state requires user
awareness.

When an email in the \textbf{Trash} is opened, a warm amber banner reads:
\textit{``This email is in Trash. Messages are permanently deleted after 30 days.''}
Two action buttons — Restore to Inbox and Delete forever — are directly embedded. The user
instantly knows the email's fate, their remaining options, and the time window.

When an email in the \textbf{Archive} is opened, a blue banner reads:
\textit{``This email is archived and not in your Inbox.''} This directly answers the
question users frequently ask in Gmail: ``where did my email go?''

After \textbf{sending}, the compose window closes and a success toast appears: ``Message
sent'' — confirming the action without requiring the user to navigate to Sent to verify.

After \textbf{login}, the inbox appears immediately and the Welcome Header greets the user
by first name with live stats — confirming successful session establishment without a
separate ``You are now logged in'' screen.


\begin{figure}[H]
  \centering
  \includegraphics[width=\linewidth]{closure.jpg}
  \caption{Trash context banner with Restore and Delete Forever actions}
\end{figure}

\begin{figure}[H]
  \centering
  \includegraphics[width=\linewidth]{archieve.jpg}
  \caption{Archive context banner: ``This email is archived and not in your Inbox'' — closure for the user about email state}
\end{figure}


% ─────────────────────────────────────────────────
\section{S5 — Permit Easy Reversal of Actions}

\begin{principleBox}{S5 · Permit Reversal}
  Users should be able to undo any reversible action.
  Reduces anxiety and encourages exploration.
\end{principleBox}

Reversal in Lucid is implemented at three levels, each matching the severity of the action:

\textbf{Level 1 — Instant undo via toast.} Archive, Trash, Move, and Mark Read operations
display a toast with an ``Undo'' button. Clicking it calls \texttt{undoLastAction()} from the
email store, which applies a stored \texttt{UndoAction} snapshot to restore the previous state.

\textbf{Level 2 — Undo keyboard shortcut.} The \texttt{Z} key triggers the same undo
mechanism from the keyboard without requiring the user to reach for the mouse. If there is
nothing to undo, a toast reads ``Nothing to undo'' — preventing a silent no-op.

\textbf{Level 3 — Soft-delete.} Deleting an email moves it to Trash, not permanent
deletion. The Trash auto-purges after 30 days. Users have three weeks to discover and recover
a misdeleted email. This matches Gmail's own model (Jakob's Law) but is made explicit with
the Trash banner (S4).

\textbf{Destructive confirmation.} Permanently deleting from Trash requires explicit
confirmation via a \texttt{ConfirmDialog} with the label ``Delete forever (cannot be undone)''.
The destructive button is rendered in red. The cancel button is the visually dominant path,
implementing the Fitts' Law principle of making the dangerous target smaller and less prominent.



\begin{figure}[H]
  \centering
  \includegraphics[width=\linewidth]{reversal.jpg}
  \caption{Trash context banner with Restore and Delete Forever actions}
\end{figure}

\begin{figure}[H]
  \centering
  \includegraphics[width=\linewidth]{forward.jpg}
  \caption{S5 — Permit Reversal: Forward action in Email Detail --- all email actions reachable and reversible}
\end{figure}

% ─────────────────────────────────────────────────
\section{S6 — Error Prevention and Handling}

\begin{principleBox}{S6 · Error Prevention}
  Prevent errors first by design; when unavoidable, provide plain-language
  recovery guidance.
\end{principleBox}

\textbf{Prevention — subject field guard.} If the user attempts to send a message with an
empty subject line, the send is intercepted. The subject field's border turns
\textcolor{danger}{red} and a warning bar appears: \textit{``Subject is empty.''} alongside
a ``Send anyway'' escape hatch — respecting autonomy (S7) while preventing the most common
composing mistake.

\textbf{Prevention — send button disablement.} The Send button is disabled when the
\texttt{To} field is empty. Users cannot send an email with no recipient, regardless of
compose window state.

\textbf{Prevention — tag name uniqueness.} In Settings → Tags, creating a duplicate tag
name triggers an immediate inline error: \textit{``A tag with this name already exists.''}
before the form is submitted — constraining the interaction to valid inputs only (Norman's
Constraints, D3).

\textbf{Prevention — email validation at login.} The login modal validates format
(\texttt{@} required) before submission and marks the field in red with
\textit{``Enter a valid email address.''} The error attaches to the specific field rather than
appearing at the top of the form — placing the error where the user's attention already is
(Raskin's Third Law).

\textbf{Recovery — compose autosave.} Every 30 seconds of inactivity, the compose window
silently saves a draft. If the browser closes or the window crashes, the draft is recoverable
from the Drafts folder — directly implementing Raskin's First Law.

\begin{figure}[H]
  \centering
  \includegraphics[width=\linewidth]{errorhand.jpg}
  \caption{Empty subject warning bar with ``Send anyway'' escape hatch}
\end{figure}
% ─────────────────────────────────────────────────
\section{S7 — Support Internal Locus of Control}

\begin{principleBox}{S7 · Locus of Control}
  The user always feels they are controlling the system, not being controlled by it.
  No auto-popups, no surprise navigation, no system-initiated changes.
\end{principleBox}

\textbf{Resizable sidebar.} The sidebar has a drag handle on its right edge. Users can pull
it to any width between 68px (compact icon mode) and 300px. At 120px it snaps to either
compact or expanded — providing predictable snap behaviour while still letting the user control
the breakpoint they want. No setting menu is required; the affordance is the handle itself.

\textbf{Settings are immediate.} In the Settings page, there is no ``Save'' button.
Every toggle, segmented control, or option card applies its change instantly. The subtitle
``Changes are saved automatically'' appears below the Signature textarea to make this explicit.
Users never have to wonder whether their preferences were committed.

\textbf{No auto-submission.} The Compose Window does not send on tab-out, on window-blur,
or on any system event. Only the Send button or \texttt{Ctrl+Enter} triggers sending. The
close button (\texttt{×}) autosaves the draft and closes — it never sends.

\textbf{Compose window controls.} The title bar exposes Minimise, Fullscreen, and Close
buttons. The window can be minimised to a 48px strip (to keep context while reading), or
expanded to a full-screen editor for long messages. The user decides the modality, not the
system.

\textbf{Account switching.} The account switcher in the Top Bar shows all signed-in
accounts. The user decides which one is active at any moment. The system never auto-switches
based on an incoming email from a different account.


\begin{figure}[H]
  \centering
  \includegraphics[width=\linewidth]{locuscontrol.jpg}
  \caption{Locus Of Control}
\end{figure}

% ─────────────────────────────────────────────────
\section{S8 — Reduce Short-Term Memory Load}

\begin{principleBox}{S8 · Reduce STM Load}
  Don't make users remember information across screens. Display what they need,
  when they need it. Related to Miller's 7$\pm$2 Law.
\end{principleBox}

\textbf{Welcome Header.} When the inbox is open, a Welcome Header above the email list
displays the active account's name, avatar, and three live statistics: Unread count,
Sent count, and Drafts count. Users never need to navigate to Sent or Drafts just to
find out how many they have.

\textbf{Unread badge in sidebar.} When the sidebar is in expanded mode, unread counts
appear as coloured badges on each nav item. In compact mode, an unread dot-badge
(\texttt{unread-dot-badge}) signals the presence of unread mail without revealing the
number — respecting the compact layout while preserving the signal.

\textbf{Recent searches.} The search dropdown surfaces the last five search queries
when focused with an empty field. Users do not need to remember what they previously
searched — the system remembers and offers recognition over recall.

\textbf{Context-aware bulk bar.} When emails are selected, the bulk action bar shows
only actions appropriate for the current folder. In Trash, it shows ``Restore'' and
``Delete forever'' — not ``Archive'' (which would be nonsensical). In Archive, it shows
``Move to Inbox'' and ``Move to Trash.'' The relevant memory of ``what folder am I in''
is embedded in the available actions, not required from the user.

\textbf{Email detail — sender and date always visible.} The sender's name, email, and
full timestamp are always displayed in the Email Detail header, regardless of scroll position.
Users who receive emails a week later do not need to remember when they arrived.


\begin{figure}[H]
  \centering
  \includegraphics[width=\linewidth]{miller.jpg}
  \caption{Reduce Short-Term Memory Load}
\end{figure}

% ══════════════════════════════════════════════════════════════
\chapter{Nielsen's 10 Usability Heuristics}
% ══════════════════════════════════════════════════════════════

% ─────────────────────────────────────────────────
\section{N1 — Visibility of System Status}

The system must always communicate its current state via timely, appropriate feedback.

In Lucid, the active navigation item receives \texttt{aria-current="page"} and is
highlighted with a blue background (\texttt{brand-50}) and bold text — making the current
location unambiguous at all times. The search bar placeholder text includes \texttt{( / )} to
signal the keyboard shortcut, surfacing system capability as status information.

The Compose Window title bar serves as a persistent status strip: it shows the email's
subject (or ``New message''), a save-status indicator (\texttt{\faCheck\ Saved} / \texttt{Saving…}),
and the window-state buttons. Even when minimised to 48px, the subject is visible — the user
always knows what message they are working on.


\begin{figure}[H]
  \centering
  \includegraphics[width=\linewidth]{inbox.jpg}
  \caption{Sidebar with ``Inbox'' active item highlighted in blue}
\end{figure}

% ─────────────────────────────────────────────────
\section{N2 — Match Between System and Real World}

Terminology and metaphors must come from the user's world, not technical internals.

Lucid avoids all internal jargon. The folders are named Inbox, Starred, Sent, Archive,
Drafts, Spam, and Trash — the same nouns English-speaking users have used in email since
the 1990s (Jakob's Law). Icons reinforce the metaphors: a trash can for Trash, an envelope
for Inbox, a star for Starred, a paper aeroplane for Send.

The compose body placeholder reads: ``Write your message… (Ctrl+B bold, Ctrl+I italic,
Ctrl+U underline)'' — plain English with inline teaching, not ``ContentEditable region
with rich text support enabled.''

Error messages use the same register: ``Enter a valid email address'' (not ``Invalid
input format''), ``Subject is empty'' (not ``Missing required header field''), ``No results
for X'' with a suggestion to ``try different keywords'' (not ``Query returned 0 records'').

% ─────────────────────────────────────────────────
\section{N3 — User Control and Freedom}

Users often take wrong turns; always provide a clear exit.

Every view in Lucid has an escape. The Email Detail has a ``Back'' button with a left arrow.
The Settings page has a back button in its header. Compose windows have a close (×) button,
a minimise button, and the keyboard \texttt{Escape} key dismisses the shortcuts overlay.
The bulk-select mode (selecting emails by checkbox) has an explicit close button (×) in the
bulk bar, returning to normal mode. No modal or mode exists in Lucid without a visible,
labelled escape.

The search bar's clear button (\texttt{×}) appears as soon as text is typed, removing the
need to manually select-all and delete — a one-click escape from the search state.


\begin{figure}[H]
  \centering
  \includegraphics[width=\linewidth]{search.jpg}
  \caption{Bulk action bar with close (×) to exit selection mode}
\end{figure}

% ─────────────────────────────────────────────────
\section{N4 — Consistency and Standards}

Users should not wonder whether different words mean the same thing.

Platform conventions are followed rigorously. The top bar layout — logo left, search
centre, account right — mirrors Google's own application chrome (Jakob's Law). Dropdown
menus all share the \texttt{.dropdown} class: same border, same border-radius, same shadow,
same item hover style. Every dropdown fades in with a \texttt{fade-up} animation at the same
speed (150ms).

Icon semantics are locked: the same star icon means starred/unstarred in the email list,
the email detail sidebar, and the bulk action bar. The same trash icon means delete in
compose (discard draft), email rows (delete), and the detail sidebar (move to trash).

% ─────────────────────────────────────────────────
\section{N5 — Error Prevention}

Design proactively eliminates error-prone situations.

Beyond the subject-guard and recipient-guard described under S6, Lucid prevents errors through
\textbf{constraint} (Norman D3). The Send button is a \texttt{disabled} HTML attribute —
not just visually dimmed — when the To field is empty, which means it is truly unreachable
by keyboard focus.

In the Tags tab, the ``Create tag'' button is the only submission path; pressing Enter
in the name field triggers the same handler, preventing accidental blank-name tags.

In the zoom widget, the \texttt{-} button is \texttt{disabled} at 70\% and the \texttt{+}
button is \texttt{disabled} at 150\%, preventing the interface from ever reaching an
unreadable state through repeated clicks.

% ─────────────────────────────────────────────────
\section{N6 — Recognition over Recall}

Show options; do not make users remember commands.

The sidebar labels display icons alongside text when expanded. In compact mode, icons alone
are shown — but every compact item has a \texttt{Tooltip} revealing its full label on hover,
bridging recognition and recall. Custom label buttons in compact mode show the label's first
letter with its unique colour — recognition of the colour replaces recall of the name.

The recipient autocomplete in the Compose Window surfaces a dropdown of matching contacts
as the user types. Users do not need to remember full email addresses — the first two
characters trigger recognition-based selection.

The search dropdown shows recent searches when the field is focused and empty — surfacing
history without requiring the user to remember what they previously typed.


\begin{figure}[H]
  \centering
  \includegraphics[width=\linewidth]{readymade_options.jpg}
  \caption{Recognition over Recall}
\end{figure}


% ─────────────────────────────────────────────────
\section{N7 — Flexibility and Efficiency of Use}

Cater to both novice and expert users.

Lucid explicitly supports two usage modes. \textbf{Novice users} navigate by clicking the
sidebar, reading tooltip labels, using the visible search bar, and clicking action buttons.
\textbf{Expert users} compose with \texttt{C}, search with \texttt{/}, archive with
\texttt{E}, navigate with \texttt{G I / G S}, and format with \texttt{Ctrl+B}.

The density setting (Compact / Default / Spacious) lets expert users who process email
quickly choose Compact — seeing more emails per screen — while less frequent users prefer
Spacious for better legibility.

The zoom widget in the Top Bar lets users with high-resolution monitors scale down to 70\%
for information density, while users with accessibility needs scale up to 150\%.

% ─────────────────────────────────────────────────
\section{N8 — Aesthetic and Minimalist Design}

Every element competes for attention. Irrelevant information reduces the visibility of important information.

Lucid's inbox view contains exactly four pieces of information per email row:
the sender avatar (identity), the subject and preview (content), the timestamp (recency),
and a star (importance). No advertisement banners, no promotional score widgets, no
``Promotions'' / ``Social'' / ``Updates'' tab noise from Gmail.

The Email Detail's action sidebar is a single column of icon-only buttons with no labels —
relying on \texttt{Tooltip} for discovery. This keeps the main reading area at full width
without a distracting toolbar row across the top.

The Settings page uses generous whitespace (40px horizontal, 48px vertical padding) and
groups each setting into a clearly titled section — consistent with Simon's Law: visual
hierarchy reduces perceived complexity.

% ─────────────────────────────────────────────────
\section{N9 — Help Users Recover from Errors}

Error messages must identify what went wrong and how to fix it.

Every error in Lucid is inline, co-located with the field that caused it, written in
first-person-familiar English. In the login modal: \textit{``Enter a valid email address''}
appears directly below the email field in \textcolor{danger}{red}. In the Tags tab:
\textit{``A tag with this name already exists.''} appears below the name input.
In the compose warning bar: \textit{``Subject is empty.''} appears immediately below the subject
row with a ``Send anyway'' escape — acknowledging the error while respecting that the user
might have a legitimate reason to send without a subject.

None of these messages reference internal states, codes, or HTTP methods.

% ─────────────────────────────────────────────────
\section{N10 — Help and Documentation}

Even self-explanatory systems need searchable, task-focused help.

The Help button (question-mark icon, top-right of the Top Bar) opens the Keyboard Shortcuts
Overlay — a task-focused reference categorised into Actions, Navigation, and Search \& Help.
The shortcut descriptions are written as verbs (``Compose new email'', ``Archive selected'',
``Focus search bar'') — task-oriented language, not feature-names.

Tooltips are the primary in-context help mechanism: every icon-only button in the sidebar,
the Top Bar, and the Email Detail sidebar has a \texttt{Tooltip} with a concise action
description, triggered on hover without interrupting the user's current task
(unobtrusiveness principle of User Support).

\begin{figure}[H]
  \centering
  \includegraphics[width=\linewidth]{shortcut_icon.jpg}
  \caption{Recognition over Recall}
\end{figure}

% ══════════════════════════════════════════════════════════════
\chapter{Norman's 6 Design Principles}
% ══════════════════════════════════════════════════════════════

Donald Norman's principles, introduced in \textit{The Design of Everyday Things},
focus on the psychology of interaction at the individual control level.

\section{D1 — Visibility}

Functions must be visible to be discoverable. Hidden equals unused.

The Compose button is the topmost element in the sidebar in both compact and expanded modes.
It is the first thing the sidebar renders, matching the Primacy Effect (L2). In compact mode
it is a 38px circle with a \texttt{+} — unmissable at any sidebar width. In expanded mode it
becomes a full-width blue pill with the word ``Compose.''

The search bar spans the entire centre of the Top Bar, with \texttt{flex: 1; max-width: 720px}.
It is always visible; it is never behind a click-to-reveal icon. This directly addresses
Gmail's mobile regression where the search field is collapsed to an icon.


\begin{figure}[H]
  \centering
  \includegraphics[width=\linewidth]{expand_closure.jpg}
  \caption{Compose button in compact (circle) and expanded (pill) sidebar modes}
\end{figure}

\section{D2 — Feedback}

Every action must produce an immediate, perceivable response.

Beyond the toast and save-status systems already described, Lucid adds micro-feedback at the
input level. The tag colour selector in Settings uses a circular button grid; the selected
colour gets a 3px dark border \textit{and} a box-shadow ring in the same colour — two
simultaneous feedback signals to accommodate both high-contrast and colour-dependent users
(Universal Design U4).

The Login button shows a spinning CSS animation (\texttt{animation: spin .7s linear infinite})
during the authentication flow — communicating ``the system is working'' without blocking
the UI or losing the user's context.

\section{D3 — Constraints}

Restrict interaction to valid actions only.

Physical constraints: the sidebar resize handle can only be dragged horizontally; the minimum
width (68px) and maximum width (300px) are enforced at the JavaScript event level.

Logical constraints: the zoom widget buttons are disabled at the boundary values. The Send
button is genuinely disabled (not just styled) when the To field is empty. The tag ``Create''
button allows submission only when the name field is non-empty.

Cultural constraints: the Delete forever button is styled in red because red universally
signals danger in Western UI conventions. The Restore button is blue (brand action) because
blue signals trusted, safe navigation.

\section{D4 — Mapping}

Controls must logically relate to their effects.

The sidebar resize handle appears on the \textit{right edge} of the sidebar — dragging
rightward makes it wider, dragging leftward makes it narrower. The direction of the drag maps
directly to the direction of the effect, matching the steering-wheel analogy (same direction).

Compose window window-control buttons appear in the \textit{title bar} of the compose window —
not in the body or the footer. This maps to the OS convention (Windows, macOS) where window
controls live at the top-right of window chrome.

The Cc and Bcc toggles in the Compose Window appear at the right edge of the To/From area —
the logical location where ``additional recipients'' would be added in a physical letter-writing
metaphor.

\section{D5 — Consistency}

Similar interface elements do similar things. Rules have no exceptions.

All left-click interactions select or activate. All hover states use the same 120ms
transition. All disabled states use \texttt{opacity: 0.5}. All icon buttons (\texttt{icon-btn})
are 36$\times$36px squares with \texttt{border-radius: var(--radius-sm)} and the same
hover background: \texttt{var(--gray-100)}.

\section{D6 — Affordance}

Object appearance communicates how it should be used.

The Compose circle button (\texttt{+}) affords clicking with a clear circular shape and a
recognisable plus symbol. The sidebar resize handle is a 4px-wide strip that changes cursor
to \texttt{col-resize} on hover — affording horizontal dragging. Every text input field has
a 1.5px border that changes colour on focus — affording text entry. The star icon affords
toggling (filled = active, outline = inactive).

The zoom widget's \texttt{-} and \texttt{+} labels directly signal their arithmetic effects
— minus decreases, plus increases — without requiring any label or documentation.


\begin{figure}[H]
  \centering
  \includegraphics[width=\linewidth]{locuscontrol.jpg}
  \caption{Sidebar resize handle with cursor affordance}
\end{figure}

% ══════════════════════════════════════════════════════════════
\chapter{Gestalt Principles}
% ══════════════════════════════════════════════════════════════

The Gestalt school of perceptual psychology provides the theoretical basis for how users
group, separate, and navigate visual information in Lucid's interface.

\section{G1 — Closure (Reification)}

The brain fills in missing information to complete familiar shapes. Lucid uses skeleton
loading rows — animated grey bars in the exact proportions of a real email row (avatar circle,
sender bar, preview bar, timestamp bar) — so that users perceive incoming data before it
arrives. The shape is recognised; the content is mentally anticipated. This prevents the
disorienting ``blank page flash'' that triggers negative visceral reactions (P3).

\section{G2 — Common Region}

Elements inside a bordered area are perceived as one group.

The bulk action bar (\texttt{bulk-bar}) is a visually distinct strip with its own background
colour (\texttt{var(--brand-50)}), top border (\texttt{1px solid var(--brand-100)}), and
height. When it appears, it is immediately perceived as a region separate from the email list
below — a ``control panel'' for the selection. Users do not confuse it with the email header.

Each email row is bounded by a subtle bottom border, which creates micro-regions that group
each email's metadata (avatar, sender, subject, preview, timestamp) into one scannable unit.

\section{G3 — Proximity}

Elements close together are perceived as related.

In the Email Detail view, the two-column layout uses physical separation as a grouping
mechanism: the action sidebar (star, archive, trash, spam) is a narrow left column, visually
close together, perceived as one group of ``actions on this email.'' The main content
(subject, sender, body, smart replies, reply/forward buttons) fills the right column —
separated by a visible border.

In the Settings page, each setting has its label and description on the left, and its
control (toggle, segmented control, option cards) on the right — maintaining proximity of
label to control throughout all six tabs.


\begin{figure}[H]
  \centering
  \includegraphics[width=\linewidth]{spam.jpg}
  \caption{Email Detail two-column layout: action sidebar (left) and content (right)}
\end{figure}

\section{G4 — Common Fate}

Elements moving together are seen as one group.

In the search dropdown, all results slide in together with the same \texttt{fade-up}
animation. The result group is perceived as a unified response to the search query,
not as individual independent items appearing sequentially.

When a compose window is minimised, the title bar collapses smoothly — the entire window
shrinks as one element, reinforcing that it is a single object changing state, not a set of
disparate controls being hidden.

\section{G5 — Symmetry}

Symmetrical compositions are perceived as stable.

The landing page hero section is centred both horizontally and vertically within the viewport.
The wordmark, headline, subheadline, CTAs, and fine print are all centre-aligned — producing
a symmetrical composition that conveys calm authority. The feature cards grid and testimonials
grid both use \texttt{auto-fit} with equal-width columns — symmetric layout signals order
and credibility.

The Compose Window is consistently right-aligned at a fixed \texttt{right: 16px} offset —
multiple windows offset by a fixed 556px — so that even with multiple compose windows open,
the layout remains geometrically predictable.

\section{G6 — Similarity}

Elements sharing visual traits are perceived as related.

All navigation items in the sidebar share the same height (40px), the same icon size
(16$\times$16px), the same font size (0.875rem), and the same hover background. They are
immediately perceived as a navigation group.

All toast notifications share the same bottom-fixed position, the same border-radius, the
same shadow, and the same typographic hierarchy — sender name bold, message body regular,
Undo link in brand colour. A user who sees one toast knows how to read every subsequent one.

\section{G7 — Figure and Ground}

The brain separates foreground objects from background.

When a compose window is open, a semi-transparent dark overlay (\texttt{rgba(32,33,36,.28)})
with a \texttt{backdropFilter: blur(3px)} is applied to the entire application behind it.
This blurs the background into the ``ground'' and brings the compose window to the
``figure'' — making the modal's priority unambiguous without completely hiding the
context behind it.

The same technique is used for the login modal on the landing page (\texttt{rgba(0,0,0,.6)}
with \texttt{blur(8px)}) — the page remains visible but subordinate.


\begin{figure}[H]
  \centering
  \includegraphics[width=\linewidth]{2.jpg}
  \caption{Email Detail two-column layout: action sidebar (left) and content (right)}
\end{figure}

\section{G8 — Continuity}

Elements aligned on a smooth path or line are seen as connected.

The Settings page left navigation is a vertical list of six tab buttons — all left-edge
aligned, all with identical spacing (2px gap). The eye travels smoothly down the list
perceiving it as a single navigation rail, not six independent controls.

In the Email Detail, the subject, sender section, and body are all left-aligned, creating a
vertical reading line that the eye follows naturally top-to-bottom without interruption.

% ══════════════════════════════════════════════════════════════
\chapter{Design Laws}
% ══════════════════════════════════════════════════════════════

\section{L1 — Miller's Law (7$\pm$2)}

Short-term memory holds approximately seven chunks of information.

The primary sidebar navigation is divided: only four items (Inbox, Starred, Sent, Archive)
are shown by default. A ``More folders'' toggle reveals the remaining four (All Mail, Drafts,
Spam, Trash). This 4+4 chunking keeps both layers within Miller's limit. The Settings page
has exactly six tabs — within the 7$\pm$2 range. The Shortcuts overlay has three categories,
each with four to six shortcuts.

\begin{figure}[H]
  \centering
  \includegraphics[width=\linewidth]{miller.jpg}
  \caption{Miller's Law}
\end{figure}

\section{L2 — Primacy and Recency Effect}

Users best remember the first and last items in a list.

In the sidebar, Inbox is always first — the most-used destination. In the login modal,
``Email \& Password'' is the first tab — the most common authentication method. In the
Settings tabs, ``Appearance'' is first (the most exploratory users try) and ``Signature''
is last (the setting users return to finish their setup). In the landing page navigation,
``Sign in'' is last (recency) and ``Get started free'' is the final call-to-action.

\begin{figure}[H]
  \centering
  \includegraphics[width=\linewidth]{recency.jpg}
  \caption{Miller's Law}
\end{figure}

\begin{figure}[H]
  \centering
  \includegraphics[width=\linewidth]{set6.jpg}
  \caption{Miller's Law}
\end{figure}

\section{L3 — Fitts' Law}

Click time increases with distance and decreases with target size.

The Compose button is the largest clickable element in the sidebar in expanded mode —
a full-width 44px-high pill — because Compose is the most frequently initiated action.
In compact mode it is a 38$\times$38px circle, meeting Apple's 44$\times$44px minimum
touch target recommendation.

The permanent-delete confirmation button in the ConfirmDialog is smaller and bordered
in red — a deliberately smaller, visually less prominent target to reduce accidental
permanent deletion.

The zoom \texttt{-} and \texttt{+} buttons are 24px squares — small but adjacent to
each other for quick fine-tuning. The spacing between them and the label is minimal,
reducing the distance users must travel between presses.


\begin{figure}[H]
  \centering
  \includegraphics[width=\linewidth]{inbox.jpg}
  \caption{Compose button as the largest target in the sidebar (Fitts' Law)}
\end{figure}

\section{L4 — Hick's Law}

Decision time increases with the number of choices.

Lucid reduces choices at every decision point. The Compose Window's formatting toolbar
has exactly seven actions (Bold, Italic, Underline, Link, Bullet list, Blockquote, Attach)
— the primary formatting operations only, not a full word-processor toolbar.

The search filters surface only two options: ``Has attachment'' (toggle) and ``From:''
(text). Advanced filter logic is not exposed at the UI level — complexity is absorbed
(Tesler's Law, L9).

The Settings tabs are labelled with single-concept nouns (Appearance, Notifications,
Shortcuts, Tags, Accounts, Signature). No tab combines two concepts, preventing the
``where would I find this?'' choice paralysis.

\begin{figure}[H]
  \centering
  \includegraphics[width=\linewidth]{compose.jpg}
  \caption{Hick's Law}
\end{figure}

\section{L5 — Power Law of Practice}

Task completion time decreases as a function of practice trials.

Lucid's keyboard shortcut system is designed to accelerate with practice. The first time
a user opens the app, they navigate via mouse. After ten sessions, the \texttt{C} key for
compose becomes automatic. The two-key sequences (\texttt{G I}, \texttt{G S}) feel
difficult on day one but are near-instant after a week — matching the logarithmic learning
curve.

The shortcut remap feature directly acknowledges that users coming from other clients
(\textit{e.g.} Superhuman: \texttt{E} for archive) have already built practice-curve
memory that should not be erased. Remapping \texttt{E} or \texttt{\#} costs zero
retraining.

\begin{figure}[H]
  \centering
  \includegraphics[width=\linewidth]{shortcut.jpg}
  \caption{Power Law of Practice}
\end{figure}

\section{L6 — Pareto Principle (80/20)}

80\% of usage comes from 20\% of features.

The primary 20\% in an email client: compose, read, search, archive, delete. Every one
of these is accessible in Lucid with a single input: \texttt{C}, click-to-read, \texttt{/},
\texttt{E}, \texttt{\#}. The remaining 80\% of features (tags, notifications, signature,
meeting links, chat) are progressively disclosed through the sidebar ``More'' toggle and
the Settings page — visible and reachable, but not competing for attention on the main
interface.

\begin{figure}[H]
  \centering
  \includegraphics[width=\linewidth]{set1.jpg}
  \caption{Pareto Principle}
\end{figure}

\section{L7 — Zipf's Principle of Least Effort}

Users want maximum outcome for minimum effort.

\textbf{Smart replies} are the clearest manifestation: after reading an email, the user
sees three one-click replies pre-generated from the email context. A single click opens
a compose window pre-filled with the response — reducing a 30-second task to 2 seconds.

\textbf{Recipient autocomplete} reduces email address entry from a full typing task to a
two-character recognition-and-select operation.

\textbf{Bulk operations} allow selecting all emails in a label with one checkbox toggle,
then archiving or deleting all with one button — compressing N repetitive actions into one.

\begin{figure}[H]
  \centering
  \includegraphics[width=\linewidth]{readymade_options.jpg}
  \caption{Smart reply buttons in Email Detail view}
\end{figure}

\section{L8 — Jakob's Law of Usability}

Users expect your site to work like every other site they use.

Lucid follows the standard web application shell:
\begin{itemize}[leftmargin=*, itemsep=2pt]
  \item Logo top-left → returns to inbox on click.
  \item Search bar centred in the top bar.
  \item Account/profile controls at top-right.
  \item Navigation in the left sidebar.
  \item Content in the main right panel.
  \item Footer on the landing page: Privacy, Terms, Security, Help.
\end{itemize}

The LUCID wordmark uses the exact Google brand colours (Blue, Red, Yellow, Green, Blue)
to leverage existing colour associations with trusted, familiar software. Users who have
used Gmail for a decade will read this as a deliberate reference — reducing the novelty
anxiety of switching clients.


\section{L9 — Tesler's Law (Conservation of Complexity)}

Every application has irreducible complexity. If removed from the UI, it moves to the developer.

The autocomplete logic in the Compose Window silently parses the To/Cc/Bcc fields, splits
on commas, isolates the last entry, and filters contacts — absorbing the complexity of
email-address parsing. The user types a name fragment and selects.

The search system applies filters from two sources (the visual filter chips and the inline
From: field) and merges them with the query string in a \texttt{useSearch} hook — the user
toggles a chip and results update; all intersection logic is invisible.

Draft autosave reconstructs a full email object (with sender, recipients, labels, thread ID)
from partial compose state every 30 seconds — the user never configures this, it just works.

\section{L10 — Raskin's 3 Laws}

\textbf{Law 1 — Never harm your work.} Autosave every 30 seconds. Soft-delete with a
30-day recovery window. Undo for all reversible actions. Permanent delete requires a
two-step confirmation.

\textbf{Law 2 — Never waste your time.} Recipient autocomplete. Smart replies. Bulk
actions. Keyboard shortcuts. The system infers the draft save ID (\texttt{draft-\{id\}})
rather than asking the user to name their drafts.

\textbf{Law 3 — Be humane.} Error messages are co-located with the field that caused them.
The subject-empty warning appears immediately below the subject row — where the user is
already looking — not at the top of the compose window where it would be outside their
locus of attention. The Caps Lock indicator example from the reference document inspired
placing all validation inline.

% ══════════════════════════════════════════════════════════════
\chapter{Visual Design Principles}
% ══════════════════════════════════════════════════════════════

\section{V1 — Unity}

All elements feel harmonious and belong together.

Lucid uses a single CSS custom-property design token system (exposed as \texttt{:root}
variables). Every colour, border-radius, shadow, and spacing value is derived from a
named token (\texttt{--brand-500}, \texttt{--radius-md}, \texttt{--shadow-xs}). This
prevents accidental inconsistency: a developer cannot accidentally use a slightly different
blue because there is only one blue in scope.

The LUCID wordmark appears in three sizes (sm, md, lg) but always uses the same five Google
colours in the same order — a unified identity across landing page, top bar, compose window
title bar, and login modal.

\begin{figure}[H]
  \centering
  \includegraphics[width=\linewidth]{inbox.jpg}
  \caption{Unity}
\end{figure}

\section{V2 — Hierarchy}

Visual weight signals importance.

In the inbox, the sender name of unread emails is set in bold; the subject and preview of
unread emails are in \texttt{var(--text-primary)} (near-black). Read emails' senders and
previews shift to \texttt{var(--text-secondary)} (grey) — a two-tier typographic hierarchy
that visually sorts the inbox by action-required state without any label or badge.

In the email detail, the subject is an \texttt{<h1>} at the largest font size (1.375rem,
bold). The sender name is secondary (0.9375rem, semi-bold). The sender email and timestamp
are tertiary (0.75rem, muted grey). Body text sits between the subject and the timestamp in
both size and visual weight — the full hierarchy is derivable from the page without
reading a word.

\begin{figure}[H]
  \centering
  \includegraphics[width=\linewidth]{preview2.jpg}
  \caption{Hierarchy}
\end{figure}

\section{V3 — Balance}

Even distribution of visual weight prevents visual unease.

The top bar achieves asymmetric balance: the logo (left) carries brand weight; the search
bar (centre) carries functional weight; the icon cluster (right) is lighter in visual mass.
Despite being different widths, the three regions feel balanced because the search bar's
\texttt{flex: 1} expansion fills the centre without over-weighting either side.

The landing page features and testimonials grids use \texttt{repeat(auto-fit,
minmax(280px, 1fr))} — auto-filling equal-width columns that maintain visual balance even
as viewport width changes.

\begin{figure}[H]
  \centering
  \includegraphics[width=\linewidth]{inbox.jpg}
  \caption{Balance : Sizing between sidebar, content, top search bar}
\end{figure}

\section{V4 — Contrast}

Differences in colour and size make elements stand out.

The primary CTA button on the landing page (``Get started free'') uses a deep blue
(\texttt{\#1A73E8}) with a glowing box-shadow (\texttt{rgba(26,115,232,.55)}) against a dark
near-black background — maximum contrast, maximum dominance. Secondary CTA (``Try the demo'')
uses a frosted glass style with 6\% white background — subordinate but visible.

In the email list, unread count badges use the brand blue on white — a contrast ratio well
above WCAG's 4.5:1 minimum. The danger-red of delete buttons against white background
provides a distinct visual channel for destructive actions.

\begin{figure}[H]
  \centering
  \includegraphics[width=\linewidth]{spam.jpg}
  \caption{Contrast}
\end{figure}

\section{V5 — Scale}

Relative sizes communicate importance.

The landing page headline uses \texttt{clamp(2.8rem, 6vw, 5rem)} — the largest text on the
page by far. The subtitle is 1.1875rem. The social proof stats (50K+, 4.9$\star$) use 2rem
bold — large enough to read at a glance. Feature card titles are 1rem, descriptions 0.9rem.
The entire size hierarchy communicates: hero claim → social proof → feature details.

\begin{figure}[H]
  \centering
  \includegraphics[width=\linewidth]{resize.jpg}
  \caption{Relative sizing}
\end{figure}

\section{V6 — Dominance}

A single element draws the eye first.

On the landing page, the LUCID wordmark rendered in Google's four-colour gradient at 3.75rem
is the dominant element of the hero. Below it, the animated twinkle of cursor-proximity stars
creates secondary visual movement that supports rather than competes with the wordmark.

In the main inbox, the Welcome Header with the user avatar and stat badges dominates the top
of the email list. Below it, the email list's consistent row height prevents any single row
from dominating — which is correct: in a list view, no single item should outrank another
unless read status (bold vs. grey) creates a temporary hierarchy.

\begin{figure}[H]
  \centering
  \includegraphics[width=\linewidth]{mainpage.jpg}
  \caption{Dominance}
\end{figure}

% ══════════════════════════════════════════════════════════════
\chapter{Universal Design Principles}
% ══════════════════════════════════════════════════════════════

Universal Design's seven principles ensure that Lucid works for users regardless of
ability, experience, language proficiency, or physical context.

\section{U1 — Equitable Use}

Every interactive element in Lucid carries both a visible label or icon and an
\texttt{aria-label} attribute for screen readers. The search input has
\texttt{role="searchbox"} and \texttt{aria-label="Search mail"}. The sidebar nav has
\texttt{role="navigation"} and \texttt{aria-label="Mail navigation"}. The email list has
\texttt{role="list"} and \texttt{aria-label="Email list"}. The compose window has
\texttt{role="dialog"} and \texttt{aria-label="Compose email"}.

No feature is available exclusively to mouse users. Every action has a keyboard pathway:
keyboard shortcuts, Tab-accessible buttons, Enter-submittable forms.

\section{U2 — Flexibility in Use}

\textcolor{red}{\textbf{Lucid accommodates right-handed, left-handed, and keyboard-only users. The sidebar canbe resized from either direction. The zoom widget scales the entire interface, including all touch targets proportionally.}}Three keyboard navigation paradigms are supported
simultaneously: single-key shortcuts, two-key navigation sequences, and modifier-key
combinations in the compose editor.

\section{U3 — Simple and Intuitive Use}

\textcolor{red}{\textbf{The login modal eliminates complexity for first-time users: three tabs (Email, Google,
GitHub) with clear labels}, a loading indicator during authentication, and a ``Continue
as Guest (Demo)'' fallback that requires zero credentials.}

Empty states across all folders use plain language: ``Your inbox is empty. When you receive
mail, it will appear here.'' There are no error codes, no ``null result set'' messages,
and no technical state descriptions exposed to the user.

\section{U4 — Perceptible Information}

\textbf{Visual + textual:} All icon buttons pair with tooltip text. Star icons are filled vs.
outlined — two visual channels for starred state (shape and fill), not just colour.

\textbf{Colour blindness:} No information is conveyed by colour alone. Unread emails are
bold (typographic channel) and \texttt{var(--text-primary)} colour (contrast channel).
Danger actions are red and labelled (``Delete forever'', ``Sign out''). The tag colour
selector uses circular buttons with border contrast rings — the selected state is indicated
by a 3px border regardless of colour.

\textbf{Motor accessibility:} All interactive elements are keyboard-reachable. The sidebar
resize handle is mouse-only, but all its effects (switching folders, reaching navigation)
have keyboard alternatives.

\section{U5 — Tolerance for Error}

\textcolor{red}{\textbf{The 30-day Trash window is the primary error-tolerance mechanism}: accidental deletion is
recoverable for four weeks. The autosave draft system ensures no compose-window content
is ever permanently lost. The ``Send anyway'' button on the empty-subject warning respects
user intent while preventing accidental sends without a subject.}

\begin{figure}[H]
  \centering
  \includegraphics[width=\linewidth]{reversal.jpg}
  \caption{Tolerance for error}
\end{figure}

\section{U6 — Low Physical Effort}

The zoom widget reduces the physical scanning distance for users who need a larger interface.
At 150\% zoom, all hit targets scale proportionally, reducing the motor precision required
for accurate clicking. Keyboard shortcuts eliminate mouse travel entirely for power users,
directly implementing this principle for expert users.

\begin{figure}[H]
  \centering
  \includegraphics[width=\linewidth]{resize.jpg}
  \caption{The zoom widget reduces the physical scanning distance for users who need a larger interface}
\end{figure}

\section{U7 — Size and Space}

All primary interactive targets meet or exceed the 44$\times$44px Apple Human Interface
Guidelines minimum. Sidebar navigation items are 40px tall. The Compose button is 38px in
diameter (compact) or 44px in expanded mode. The Account Pill button is 36px tall. The zoom
widget's \texttt{-}/\texttt{+} buttons are 24$\times$24px — below the 44px ideal, but
grouped together within a larger 80px interactive zone that encompasses both and the label.

% ══════════════════════════════════════════════════════════════
\chapter{Web Usability Principles (Krug)}
% ══════════════════════════════════════════════════════════════

\section{W1 — Don't Make Users Think}

Every navigation item has a self-evident label. The search bar placeholder includes the
keyboard shortcut. The Compose button says ``Compose.'' The Archive action says ``Archive.''
The permanent delete confirm dialog says ``Delete forever (cannot be undone)'' — there is
zero ambiguity about the consequences.

The F-pattern reading axis (top-left dominance) is respected: the most critical action
(Compose) is at the top of the sidebar; the most critical navigation (Inbox) is immediately
below it; the email list starts from the top of the content area; the email preview
immediately shows the sender, subject, and preview text.

\begin{figure}[H]
  \centering
  \includegraphics[width=\linewidth]{mainpage.jpg}
  \caption{Self-Evident Navigation Labels}
\end{figure}
\section{W2 — Don't Squander Users' Patience}

The landing page's ``Try the demo'' button bypasses authentication entirely — a guest login
delivers the full inbox experience in one click, without asking for an email address. The
full login flow is available, but never mandatory for initial exploration.

\begin{figure}[H]
  \centering
  \includegraphics[width=\linewidth]{freedemo.jpg}
  \caption{One-Click Guest Demo Access --- ``Try the demo'' bypasses authentication entirely}
\end{figure}
\section{W3 — Manage Users' Attention}

The landing page uses bold typographic hierarchy, a dynamic star animation, and a browser
mockup preview to guide attention top-to-bottom. Static layout above (wordmark, headline)
draws first attention; animated stars add subtle dynamism without competing with the text.

In the main application, attention management is the inverse: the interface is deliberately
static. There are no pulsing badges, no auto-playing video, no moving content that would
distract from the act of reading email.

\begin{figure}[H]
  \centering
  \includegraphics[width=\linewidth]{mainpage.jpg}
  \caption{Attention Hierarchy on the Landing Page}
\end{figure}
\section{W4 — Strive for Feature Exposure}

The four most-used features — search, compose, primary navigation, and account switching —
are all visible without scrolling or clicking any reveal button. The search bar and compose
button are above the fold at every viewport width. The sidebar primary nav (Inbox, Starred,
Sent, Archive) is always visible in compact mode.

\begin{figure}[H]
  \centering
  \includegraphics[width=\linewidth]{inbox.jpg}
  \caption{Primary Features Visible Without Scrolling --- search, compose, nav, account all above the fold}
\end{figure}
\section{W5 — Make Use of Effective Writing}

Every toast message is five words or fewer: ``Archived'', ``Moved to Trash'', ``Message sent'',
``Restored to Inbox.'' Tooltip labels are action verbs: ``Compose new email'', ``Go to inbox'',
``Clear search.'' The Settings subtitles are objective descriptions, not marketing copy:
``Controls how much vertical space each email row occupies.''

\begin{figure}[H]
  \centering
  \includegraphics[width=\linewidth]{feedback2.jpg}
  \caption{Concise Toast Messages --- five words or fewer, plain-English action confirmation}
\end{figure}
\section{W6 — Strive for Simplicity}

The KISS principle governs feature scope. The compose toolbar has seven formatting tools, not
thirty. The sidebar nav has eight folders, not dozens of configurable views. Settings has
six tabs addressing the six most common configuration needs. Nothing ships without a usability
justification — which the landing page explicitly states as the ``Lucid promise.''

\begin{figure}[H]
  \centering
  \includegraphics[width=\linewidth]{compose.jpg}
  \caption{Simplified Compose Toolbar --- seven formatting actions, not a full word-processor toolbar}
\end{figure}
\section{W7 — Don't Be Afraid of White Space}

The Settings page uses 48px horizontal padding and 40px vertical section spacing.
The Email Detail uses 24px of bottom margin after the action buttons before the reply area.
The landing page features section uses 60px above the card grid and 20px between cards.
White space in Lucid is not absence of content — it is the primary mechanism for grouping,
separation, and cognitive breathing room.

\begin{figure}[H]
  \centering
  \includegraphics[width=\linewidth]{set4.jpg}
  \caption{Generous White Space in Settings --- 48px horizontal padding, 40px vertical section spacing}
\end{figure}
\section{W8 — Communicate with Visible Language}

\textbf{Organise:} Consistent grid, consistent nav positions, consistent component shapes.

\textbf{Economise:} Maximum two typefaces (Inter/system-ui sans-serif + monospace for
shortcut key badges). Maximum three font sizes per screen. Maximum 80 characters per
email preview line.

\textbf{Communicate:} All interactive elements are visually distinct from static elements.
All error states use red. All success states use green. All brand actions use blue.

\begin{figure}[H]
  \centering
  \includegraphics[width=\linewidth]{inbox.jpg}
  \caption{Visual Language --- blue active state, bold unread, muted read; colour-coded throughout}
\end{figure}
\section{W9 — Conventions Are Our Friends}

The top bar layout follows the Google Material / Gmail convention. The sidebar-left,
content-right layout follows every established email client since 2000. The floating action
button (Compose circle) is at the top of the nav, matching both Gmail mobile and desktop.
Attachment previews show filename and file size in parentheses — matching OS file browser
conventions.

\begin{figure}[H]
  \centering
  \includegraphics[width=\linewidth]{mainpage3.jpg}
  \caption{Familiar Layout Conventions --- logo left, search centre, account right, following Google Material convention}
\end{figure}
\section{W10 — Test Early, Test Often}

The landing page's ``No credit card required · Free forever plan'' copy is a direct
usability principle: remove friction from the evaluation path so that testing is possible
for any user, at any point. The guest demo bypasses authentication for zero-friction testing.

\begin{figure}[H]
  \centering
  \includegraphics[width=\linewidth]{guest.jpg}
  \caption{Zero-Friction Guest Demo for Testing --- no credit card, no signup required}
\end{figure}
\begin{figure}[H]
  \centering
  \includegraphics[width=\linewidth]{set6.jpg}
  \caption{Settings page: generous whitespace and organised sections --- visible language, economy, and communication}
\end{figure}

% ══════════════════════════════════════════════════════════════
\chapter{Psychological Principles}
% ══════════════════════════════════════════════════════════════

\section{P1 — Aesthetic Usability Effect}

Attractive designs are perceived as more usable, even if they aren't.

The landing page leads with the most visually ambitious element in the project: a
\texttt{<canvas>}-based animated night sky with 280 twinkle-animated stars that respond
to cursor proximity with a radial blue glow. This is not functional — it does not send
email. But it establishes a visceral first impression (P3) of craftsmanship and care that
directly increases the perceived usability and trustworthiness of the entire application.

The LUCID wordmark in Google colours, the frosted-glass login modal with
\texttt{border-radius: 24px} and \texttt{box-shadow: 0 32px 80px rgba(0,0,0,.25)}, and the
browser chrome mockup in the hero section all signal: this product is finished, polished,
and professional.

\begin{figure}[H]
  \centering
  \includegraphics[width=\linewidth]{mainpage.jpg}
  \caption{P1 — Aesthetic Usability: polished landing page builds perceived trustworthiness before the user interacts}
\end{figure}

\section{P2 — Von Restorff Effect (Isolation Effect)}

Distinctive items are remembered better. One different element stands out.

Lucid's landing page has a single primary CTA — ``Get started free'' in brand blue.
Everything around it is either dark background (hero) or white card (feature section) or
dark testimonial cards. The blue button is the only blue element on the page — it cannot
be missed.

In the inbox, the only coloured element on a typical email row is the avatar (a coloured
circle unique to each sender). All rows share the same white background, same grey text,
same layout — so when one row is blue-highlighted (currently previewed email), it is
immediately distinctive.

\begin{figure}[H]
  \centering
  \includegraphics[width=\linewidth]{starred.jpg}
  \caption{P2 — Von Restorff: starred emails as a single distinctive element amid uniform rows}
\end{figure}

\section{P3 — Visceral Reactions}

Users form instant emotional judgments about design.

The night sky landing page is specifically designed to trigger a ``wow'' reaction within
the first 100ms of the page load — before the user reads a single word. The star proximity
interaction (stars brightening when the cursor is near) rewards mouse movement with a
haptic-like visual response, creating an immediate sense of a living, reactive interface.

The LUCID wordmark colour gradient (Google blue-red-yellow-green) triggers an emotional
association with a trusted, familiar brand ecosystem, reducing the ``I've never heard of
this'' scepticism that new applications must overcome.

\begin{figure}[H]
  \centering
  \includegraphics[width=\linewidth]{mainpage3.jpg}
  \caption{P3 — Visceral Reaction: night-sky canvas and Google-coloured wordmark trigger instant positive emotional response}
\end{figure}

\section{P4 — Dual-Coding Theory}

Words and images together produce stronger memory encoding than either alone.

Each feature card on the landing page pairs an emoji icon (visual channel) with a title and
description (verbal channel): \texttt{\faBolt\ Blazing fast}, \texttt{\faMoon\ Night \& day},
\texttt{\faReply\ Undo anything}. This emoji+text pairing is not decorative — it creates a
dual-coded memory trace: users who forget the title may still remember the lightning bolt.

Smart reply buttons (``Sounds great!'', ``Thanks for letting me know.'') pair an
implicit tone signal (the choice of words) with the visual affordance of a pill button — the
visual form signals ``this is a complete, ready-to-use response,'' the words signal the
emotional register.

\begin{figure}[H]
  \centering
  \includegraphics[width=\linewidth]{readymade_options.jpg}
  \caption{P4 — Dual-Coding: smart reply pill buttons pair visual affordance (shape) with verbal tone (words)}
\end{figure}

\section{P5 — Cialdini's 6 Principles of Persuasion}

The landing page deploys all six principles systematically:

\begin{tabularx}{\textwidth}{lX}
  \toprule
  \textbf{Principle} & \textbf{Implementation in Lucid}\\
  \midrule
  \textbf{Reciprocity} & ``Free forever plan'', ``No credit card required'' — giving value before asking for commitment\\
  \textbf{Authority} & ``Built on 40 years of HCI research'', ``Research-backed'' badge on the HCI-designed card\\
  \textbf{Liking} & Relatable brand voice (``Your inbox deserves better.''), first-name personalisation after login\\
  \textbf{Social Proof} & ``50,000+ users'', ``4.9\faStar\ User rating'', three named testimonials with roles (Figma designer, Atlassian engineer, Google UX researcher)\\
  \textbf{Scarcity} & ``Early access'' badge with pulsing blue dot — creates time-sensitive membership perception\\
  \textbf{Commitment} & ``Try the demo'' low-commitment entry before account creation; once using the product, users are invested\\
  \bottomrule
\end{tabularx}

\begin{figure}[H]
  \centering
  \includegraphics[width=\linewidth]{freedemo.jpg}
  \caption{P5 — Cialdini: social proof stats, ``Early access'' scarcity badge, and zero-commitment demo entry on the landing page}
\end{figure}

\section{P6 — Mental Models}

Users' beliefs about how a system works, formed from prior experience, are honoured throughout Lucid.

Gmail's mental model — sidebar-left, email-list-right, compose-as-floating-modal — is
preserved. The two-key shortcut sequences (\texttt{G I}, \texttt{G S}) are borrowed
directly from Gmail's own shortcut system so that Gmail power users face zero relearning.

The LUCID wordmark's Google colours are a deliberate mental model transfer: the product
looks like a Google product, inheriting the reliability and familiarity associations.

Where Lucid deliberately \textit{breaks} Gmail's mental model — for example, by showing a
Welcome Header with stats instead of just starting the email list — a clear label
(``Welcome, [Name]'') provides a transition anchor so users understand what they are
seeing without disorientation.

\begin{figure}[H]
  \centering
  \includegraphics[width=\linewidth]{inbox.jpg}
  \caption{P6 — Mental Models: sidebar-left, list-right layout matches Gmail's model exactly to eliminate relearning}
\end{figure}

\section{P7 — Skeuomorphism}

Interface elements mimic real-world counterparts to reduce cognitive load.

The attachment display uses a paperclip icon alongside the filename and file size —
mirroring the physical metaphor of a paper attachment to a letter. The email body uses
\texttt{blockquote} formatting (visually indented with a left border) for quoted text in
replies — matching the typographic convention of physical quoted letters.

The compose window uses \texttt{<label>} elements for ``To'', ``Cc'', ``Bcc'' —
mirroring the physical envelope fields users have seen since postal mail. The floating
window itself mimics a desktop application window with a title bar, window controls, and
a resizable body.

\begin{figure}[H]
  \centering
  \includegraphics[width=\linewidth]{compose.jpg}
  \caption{P7 — Skeuomorphism: compose window with To/Cc/Bcc envelope fields mirrors physical letter-writing metaphor}
\end{figure}

\begin{figure}[H]
  \centering
  \includegraphics[width=\linewidth]{reply_email.jpg}
  \caption{P7 — Skeuomorphism: reply compose with blockquote-indented original text mimics physical letter quoting convention}
\end{figure}

\section{P8 — Colour Psychology}

\begin{tabularx}{\textwidth}{lX}
  \toprule
  \textbf{Colour} & \textbf{Use in Lucid}\\
  \midrule
  \textbf{Blue} & Primary brand actions, active nav items, primary CTA buttons, unread count badges — conveys trust, calm, reliability\\
  \textbf{Red} & Danger actions only (Delete, Delete forever, Sign out) — universally signals stop/danger\\
  \textbf{Green} & Success states (Saved indicator, Restore to Inbox toast) — signals safety and completion\\
  \textbf{Amber/Yellow} & Trash context banner, warning — signals caution without the urgency of red\\
  \textbf{Dark navy} & Landing page background, compose window title bar — conveys sophistication and focus\\
  \textbf{White} & Main application surfaces — purity, simplicity, readability\\
  \bottomrule
\end{tabularx}

\begin{figure}[H]
  \centering
  \includegraphics[width=\linewidth]{preview1.jpg}
  \caption{P8 — Colour Psychology: amber Trash banner (caution), red Delete button (danger), green Restore (safe action)}
\end{figure}

% ══════════════════════════════════════════════════════════════
\chapter{Interaction Styles}
% ══════════════════════════════════════════════════════════════

Lucid deliberately combines multiple interaction styles to serve both novice and expert users.

\section{Form Fill-In (I2)}

The login modal, compose window (To, Cc, Bcc, Subject fields), tag creation form, and
signature textarea all use form fill-in interaction. Fields have visible labels, logical
tab order, and inline validation. The compose window uses progressive disclosure: Cc and Bcc
rows are hidden until their toggle buttons are clicked, reducing cognitive load for users
who only need ``To''.

\begin{figure}[H]
  \centering
  \includegraphics[width=\linewidth]{login_email.jpg}
  \caption{Form Fill-In (I2): login modal with labelled fields, tab order, and inline validation}
\end{figure}

\section{Menu Selection (I3)}

The sidebar, Settings navigation, and the account switcher dropdown are menu selection
interfaces. They present visible options with consistent item styles, following the 7$\pm$2
rule. The search results dropdown is a recognise-don't-recall menu: the user sees matching
emails and selects, rather than typing a full query from memory.

\begin{figure}[H]
  \centering
  \includegraphics[width=\linewidth]{search.jpg}
  \caption{Menu Selection (I3): search results dropdown --- recognise and select rather than recall and type}
\end{figure}

\section{Direct Manipulation (I4)}

The sidebar resize handle is the clearest example of direct manipulation in Lucid. The
user acts directly on a visible object (the edge of the sidebar) with an immediately
reversible action (dragging), and receives immediate visual feedback (the sidebar width
changes as the mouse moves). No command syntax is required.

Star toggling on email rows (click the star → it fills → email is starred), checkbox
selection in the email list, and drag-resizable compose window all follow the same WYSIWYG
model: visible objects, rapid reversible actions, immediate feedback.

\begin{figure}[H]
  \centering
  \includegraphics[width=\linewidth]{locuscontrol.jpg}
  \caption{Direct Manipulation (I4): sidebar drag-resize handle --- the user acts on a visible object with immediate visual feedback}
\end{figure}

\section{Natural Language (I5)}

The search bar accepts free-form natural language queries. The useSearch hook filters by
matching the query against sender name, sender email, subject, and body preview — letting
users search as they would speak: ``design meeting'', ``Priya invoice'', ``report Q3''.
The system handles the parsing; the user uses their own language.

Smart replies are pre-generated natural language responses — the system infers tone
and proposes language, which the user confirms or edits. This is a natural language
\textit{output} interaction, not input, which avoids the ambiguity risk of NL input.

\begin{figure}[H]
  \centering
  \includegraphics[width=\linewidth]{full_email.jpg}
  \caption{Natural Language (I5): smart replies as pre-generated NL output --- the user confirms tone without composing from scratch}
\end{figure}

\begin{figure}[H]
  \centering
  \includegraphics[width=\linewidth]{email_options.jpg}
  \caption{Natural Language \& Direct Manipulation: email action options accessible from the email detail view}
\end{figure}

% ══════════════════════════════════════════════════════════════
\chapter{How Lucid Fixes Gmail's Known Usability Issues}
% ══════════════════════════════════════════════════════════════

This chapter systematically addresses documented usability complaints in Gmail, mapping each
issue to a specific Lucid design decision and the underlying HCI principle.

\begin{figure}[H]
  \centering
  \includegraphics[width=\linewidth]{inbox.jpg}
  \caption{Lucid inbox: single unified inbox, no tab fragmentation, all emails in chronological order}
\end{figure}

%─────────────────────────────────────────────────
\section{Issue 1 — Gmail's Tab System Buries Important Emails}

\begin{gmailIssueBox}{Gmail Issue: Promotions / Social / Updates Tabs}
  Gmail's five-tab inbox (Primary, Social, Promotions, Updates, Forums) fragments the
  user's attention and routinely buries important emails in non-Primary tabs. Users
  frequently miss time-sensitive messages.
\end{gmailIssueBox}

\begin{lucidSolutionBox}{Lucid Solution: Single Unified Inbox with Custom Tags}
  Lucid presents a single Inbox with all messages in chronological order. Users who
  want categorisation create \textbf{custom Tags} (Settings → Tags) with their own
  colours and names — not a pre-imposed, opaque categorisation algorithm.
\end{lucidSolutionBox}

\textbf{Principle citations:} S7 (Locus of Control — user controls organisation),
N7 (Flexibility — power users create tags; casual users don't),
L9 (Tesler's Law — categorisation complexity moved to user choice, not system imposition).

%─────────────────────────────────────────────────
\section{Issue 2 — Gmail's Cognitive Load from Feature Density}

\begin{gmailIssueBox}{Gmail Issue: Overwhelming Feature Surface Area}
  Gmail exposes Meet links, Chat, Tasks, Calendar events, Google Workspace integrations,
  Smart Compose suggestions, and multiple toolbars simultaneously — creating a high
  cognitive load for users who simply want to read and write email.
\end{gmailIssueBox}

\begin{lucidSolutionBox}{Lucid Solution: Progressive Disclosure by Section}
  The sidebar divides features into clearly labelled sections (Menu / Meet / Chat) with
  a section divider and label. The four primary nav items (Inbox, Starred, Sent, Archive)
  are always visible; secondary items (Drafts, Spam, Trash, All Mail) are behind a
  ``More folders'' toggle.
\end{lucidSolutionBox}

\textbf{Principle citations:} L4 (Hick's Law — fewer simultaneous choices),
N8 (Aesthetic Minimalism — every element competes for attention),
W6 (Simplicity — KISS principle).

\begin{figure}[H]
  \centering
  \includegraphics[width=\linewidth]{personal.jpg}
  \caption{Issue 2 — Lucid sidebar: clearly labelled Menu / Meet / Chat sections; secondary folders hidden behind ``More''}
\end{figure}

\begin{figure}[H]
  \centering
  \includegraphics[width=\linewidth]{chat.jpg}
  \caption{Issue 2 — Progressive Disclosure: Chat contacts accessible in sidebar without cluttering the primary email navigation}
\end{figure}

\begin{figure}[H]
  \centering
  \includegraphics[width=\linewidth]{gmeet.jpg}
  \caption{Issue 2 — Progressive Disclosure: Google Meet link section revealed in sidebar without imposing on the primary email flow}
\end{figure}

%─────────────────────────────────────────────────
\section{Issue 3 — Gmail's Lack of Visible Undo for Archive and Delete}

\begin{gmailIssueBox}{Gmail Issue: Undo Send Only — No Undo for Archive or Delete}
  Gmail offers ``Undo Send'' within a configurable window (up to 30 seconds), but
  archiving or deleting emails provides only a brief, easily missed yellow toast
  with an ``Undo'' link that many users never notice in time.
\end{gmailIssueBox}

\begin{lucidSolutionBox}{Lucid Solution: Persistent Undo Toast + Keyboard Undo}
  Every reversible action in Lucid produces a toast with an explicit ``Undo'' button
  (not just a link). Additionally, the keyboard shortcut \texttt{Z} triggers
  \texttt{undoLastAction()} at any time without needing to catch the toast.
  The undo system uses a typed \texttt{UndoAction} snapshot that can restore the
  exact previous state.
\end{lucidSolutionBox}

\textbf{Principle citations:} S5 (Permit Reversal), N3 (User Control and Freedom),
L10 (Raskin's First Law — never allow work to be harmed), G6 (Similarity — toasts always
in the same style and position, making them recognisable).

\begin{figure}[H]
  \centering
  \includegraphics[width=\linewidth]{feedback3.jpg}
  \caption{Issue 3 — Persistent Undo Toast: explicit ``Undo'' button on every reversible action, not a fleeting link}
\end{figure}

%─────────────────────────────────────────────────
\section{Issue 4 — Gmail's Password Field Always Masked}

\begin{gmailIssueBox}{Gmail Issue: Mandatory Password Masking}
  Gmail always masks the password field with dots, which Nielsen Norman Group research
  shows reduces usability without meaningfully improving security for private-device logins.
  Users make more typos, use simpler passwords, and lose confidence in their entries.
\end{gmailIssueBox}

\begin{lucidSolutionBox}{Lucid Solution: Password Shown by Default}
  The login modal's password field uses \texttt{type="text"} with a label
  ``Password (visible)'' — showing the password in plaintext by default. The code
  comment explicitly cites NNG's password masking research as justification.
\end{lucidSolutionBox}

\textbf{Principle citations:} N5 (Error Prevention — fewer typos mean fewer login
failures), W6 (Simplicity — remove unnecessary security theatre), NNG password masking
research (web usability evidence).

\begin{figure}[H]
  \centering
  \includegraphics[width=\linewidth]{login_email.jpg}
  \caption{Issue 4 — Password Shown by Default: plaintext password field reduces login errors without security risk on private devices}
\end{figure}

%─────────────────────────────────────────────────
\section{Issue 5 — Gmail's Non-Resizable Sidebar}

\begin{gmailIssueBox}{Gmail Issue: Fixed Sidebar Width}
  Gmail's sidebar cannot be resized. Users with many custom labels have to scroll the
  sidebar. Users who want more reading space for email content cannot reclaim it from
  the sidebar.
\end{gmailIssueBox}

\begin{lucidSolutionBox}{Lucid Solution: Drag-Resizable Sidebar with Snap Behaviour}
  The sidebar has a drag handle that allows widths from 68px to 300px. Widths below
  120px snap to the 68px compact mode (icon-only); widths above 120px snap to 220px
  expanded mode. Custom snap thresholds prevent the sidebar from getting stuck in an
  unusable intermediate state.
\end{lucidSolutionBox}

\textbf{Principle citations:} S7 (Locus of Control — user controls layout),
D4 (Mapping — drag direction maps to width change),
D6 (Affordance — resize cursor signals draggability),
U2 (Flexibility in Use).

\begin{figure}[H]
  \centering
  \includegraphics[width=\linewidth]{resize.jpg}
  \caption{Issue 5 — Drag-Resizable Sidebar: compact icon mode (68px) vs. expanded text mode (220px) with snap thresholds}
\end{figure}

%─────────────────────────────────────────────────
\section{Issue 6 — Gmail's Search Has No Persistent Suggestion History}

\begin{gmailIssueBox}{Gmail Issue: Ephemeral Search Suggestions}
  Gmail's search suggestions disappear after clearing the field. Users who frequently
  search for the same queries (a colleague's name, a project keyword) must retype the
  full query each time.
\end{gmailIssueBox}

\begin{lucidSolutionBox}{Lucid Solution: Persistent Recent Searches via localStorage}
  Lucid stores up to five recent searches in \texttt{localStorage} under
  \texttt{hcimail-recent-searches}. Focusing the search bar with an empty field
  immediately surfaces these searches as clickable items — zero characters needed to
  re-execute a previous query.
\end{lucidSolutionBox}

\textbf{Principle citations:} S8 (Reduce STM Load — the system remembers so the
user doesn't have to), N6 (Recognition over Recall — user selects from history rather
than recalling the exact query), L7 (Zipf's Least Effort — most effort for repeated
tasks is reduced to zero).

\begin{figure}[H]
  \centering
  \includegraphics[width=\linewidth]{search.jpg}
  \caption{Issue 6 --- Persistent Recent Searches: focusing the empty search bar surfaces the last five queries from \texttt{localStorage}}
\end{figure}

%─────────────────────────────────────────────────
\section{Issue 7 — Gmail Keyboard Shortcuts Are Not Customisable}

\begin{gmailIssueBox}{Gmail Issue: Fixed, Non-Remappable Shortcuts}
  Gmail's keyboard shortcuts are fixed. Users who switch from other clients
  (Superhuman, Fastmail) must retrain existing muscle memory. There is no way to
  remap even frequently used actions.
\end{gmailIssueBox}

\begin{lucidSolutionBox}{Lucid Solution: Per-Shortcut Remapping in Settings}
  Every shortcut in Lucid can be remapped through Settings → Shortcuts by clicking
  a key badge and pressing a new key. Changes are applied immediately without a page
  reload. The \texttt{customShortcuts} record in \texttt{userPreferences} overrides the
  default for that shortcut ID throughout the entire keyboard handler.
\end{lucidSolutionBox}

\textbf{Principle citations:} S2 (Enable Shortcuts), N7 (Flexibility and Efficiency),
L5 (Power Law of Practice — prevent users from having to unlearn hard-won muscle memory),
S7 (Locus of Control — user controls the tool's behaviour).

\begin{figure}[H]
  \centering
  \includegraphics[width=\linewidth]{set3.jpg}
  \caption{Issue 7 — Per-Shortcut Remapping: Settings → Shortcuts tab with clickable key badges for every action}
\end{figure}

%─────────────────────────────────────────────────
\section{Issue 8 — Gmail's Compose Window Is Not Scalable}

\begin{gmailIssueBox}{Gmail Issue: Fixed-Size Compose Window}
  Gmail's compose window has three sizes: normal (popup), full-screen, and minimised.
  The normal popup size is fixed — users cannot resize it, and the body textarea is often
  too small for long emails.
\end{gmailIssueBox}

\begin{lucidSolutionBox}{Lucid Solution: Resizable Textarea + Fullscreen Mode}
  The Compose Window body uses \texttt{resize: vertical} on the textarea, letting users
  manually expand the editing area. The Fullscreen toggle (\texttt{\faExpand}) expands the
  compose window to \texttt{position: fixed; inset: 16px} — a large full-viewport editor
  ideal for composing long-form content. Minimise collapses to a 48px title strip to
  free the screen while preserving compose context.
\end{lucidSolutionBox}

\textbf{Principle citations:} S7 (Locus of Control), U2 (Flexibility in Use),
N7 (Flexibility and Efficiency), D4 (Mapping — the \faExpand\ icon signals expansion direction).

\begin{figure}[H]
  \centering
  \includegraphics[width=\linewidth]{compose.jpg}
  \caption{Issue 8 — Resizable Compose Window: vertically resizable textarea and fullscreen toggle for long-form editing}
\end{figure}

%─────────────────────────────────────────────────
\section{Issue 9 — Gmail Does Not Show Context When Viewing Archived/Trashed Email}

\begin{gmailIssueBox}{Gmail Issue: Silent Archive/Trash State}
  When opening an email from Trash or Archive in Gmail, there is no prominent visual
  signal of its state. Users who clicked through from a search result must remember
  (short-term memory load) that the email is not in their Inbox and act accordingly.
\end{gmailIssueBox}

\begin{lucidSolutionBox}{Lucid Solution: Prominent Context Banners}
  Lucid renders a full-width coloured context banner at the top of the Email Detail
  when viewing Trash (amber: ``This email is in Trash. Messages are permanently deleted
  after 30 days.'') or Archive (blue: ``This email is archived and not in your Inbox.'').
  Both banners include action buttons (Restore, Delete forever) for immediate decision-making.
\end{lucidSolutionBox}

\textbf{Principle citations:} S4 (Dialogs for Closure), S8 (Reduce STM Load — state
is displayed, not remembered), N1 (Visibility of System Status), G2 (Common Region —
the banner is a visually bounded information region).

\begin{figure}[H]
  \centering
  \includegraphics[width=\linewidth]{closure.jpg}
  \caption{Issue 9 — Trash context banner: amber warning with Restore and Delete Forever actions embedded}
\end{figure}

%─────────────────────────────────────────────────
\section{Issue 10 — Gmail Has No Interface Zoom Control}

\begin{gmailIssueBox}{Gmail Issue: No Built-in UI Scaling}
  Gmail relies entirely on the browser's text-size setting for accessibility scaling,
  which affects the entire browser and often breaks third-party page layouts. There is
  no per-app zoom control.
\end{gmailIssueBox}

\begin{lucidSolutionBox}{Lucid Solution: App-Level Zoom Widget + Settings Options}
  The Top Bar includes a \texttt{-/+} zoom widget that scales the entire UI by setting
  \texttt{document.documentElement.style.fontSize} proportionally (70\%–150\% of the
  14px base size). Settings → Appearance offers three option cards (85\%, 100\%, 115\%)
  for one-click preset selection. Because the entire application uses relative units
  (\texttt{rem}, \texttt{em}) throughout, scaling the root font size scales everything
  proportionally.
\end{lucidSolutionBox}

\textbf{Principle citations:} U2 (Flexibility in Use), U7 (Size and Space),
S7 (Locus of Control), N7 (Flexibility and Efficiency), Universal Design accessibility
compliance.

\begin{figure}[H]
  \centering
  \includegraphics[width=\linewidth]{resize.jpg}
  \caption{Issue 10 --- App-Level Zoom Widget: Appearance settings with 85\%/100\%/115\% preset cards and top-bar ±~widget}
\end{figure}

\begin{figure}[H]
  \centering
  \includegraphics[width=\linewidth]{set1.jpg}
  \caption{Issue 10 --- App-Level Zoom in Settings}
\end{figure}

%─────────────────────────────────────────────────
\section{Issue 11 — Gmail's Empty States Are Generic}

\begin{gmailIssueBox}{Gmail Issue: Uniform Empty State}
  Gmail uses a generic ``No conversations'' message across all folders. A user
  who accidentally archives all their email sees the same message as a user who
  genuinely has an empty inbox — no context, no guidance, no action.
\end{gmailIssueBox}

\begin{lucidSolutionBox}{Lucid Solution: Context-Specific Empty States}
  Lucid defines a unique empty state configuration for every folder:
  Inbox (``Your inbox is empty. When you receive mail, it will appear here.''),
  Sent (``Nothing sent yet.'' + Compose button), Drafts (``No drafts.'' + Start
  composing button), Trash (``Trash is empty. Deleted messages appear here and are
  permanently removed after 30 days.''), etc. Search also has its own distinct state
  with an illustrated icon.
\end{lucidSolutionBox}

\textbf{Principle citations:} N2 (Match System to Real World — folder-specific language),
N9 (Help Recover from Errors — guidance on what to do next), S4 (Closure — users
understand the current state completely).

\begin{figure}[H]
  \centering
  \includegraphics[width=\linewidth]{sent.jpg}
  \caption{Issue 11 --- Context-Specific Empty State: Sent folder with folder-specific guidance and Compose CTA}
\end{figure}

%─────────────────────────────────────────────────
\section{Issue 12 — Gmail Has No Welcome/Stats Header}

\begin{gmailIssueBox}{Gmail Issue: No At-a-Glance Stats}
  Gmail's inbox header shows only the folder name. To know how many unread emails,
  sent emails, or drafts exist, users must navigate to each folder individually.
\end{gmailIssueBox}

\begin{lucidSolutionBox}{Lucid Solution: Welcome Header with Live Stats}
  Lucid's inbox view opens with a Welcome Header showing the user's avatar, first name,
  custom tag chips, and three live counters: Unread, Sent, Drafts — computed from
  the live email store. No navigation required.
\end{lucidSolutionBox}

\textbf{Principle citations:} S8 (Reduce STM Load), N1 (Visibility of System Status),
V6 (Dominance — the header draws the eye to the most important contextual information
before the email list begins).

\begin{figure}[H]
  \centering
  \includegraphics[width=\linewidth]{welcome.jpg}
  \caption{Issue 12 --- Welcome Header with Live Stats: unread, sent, and draft counters visible without folder navigation}
\end{figure}

%─────────────────────────────────────────────────
\section{Issue 13 — Gmail's Reload Gives No Feedback or Closure}

\begin{gmailIssueBox}{Gmail Issue: Silent Reload — No Status Signal}
  Gmail's reload button (the circular-arrow icon next to the select-all checkbox)
  produces no loading indicator on click. On a first press it sometimes shows a
  one-frame spinner, but gives no closure message on completion. Subsequent presses
  produce nothing at all. Users are left uncertain whether the reload occurred,
  whether any new mail arrived, and whether the system is still working.
  This is a textbook violation of Nielsen's Heuristic 1 (Visibility of System Status).
\end{gmailIssueBox}

\begin{lucidSolutionBox}{Lucid Solution: Three-Phase Reload with Full Closure}
  Lucid's reload is designed with three distinct HCI-grounded phases:
  \begin{enumerate}[leftmargin=*, itemsep=2pt]
    \item \textbf{Immediate icon feedback (S3):} The \texttt{IconRefresh} circular-arrow
          icon begins spinning at 0.75s/revolution the moment the button is clicked,
          signalling that work has begun — zero ambiguity about whether the press registered.
    \item \textbf{Skeleton loading rows (G1 Closure):} The email list is replaced by
          eight animated skeleton rows in the exact proportions of real email rows
          (avatar circle, sender bar, preview bar, timestamp bar). The brain perceives
          the incoming content before it arrives, eliminating the disorienting blank-page
          flash.
    \item \textbf{Closure toast on completion (S4):} When the reload finishes, a
          contextual toast fires — \textit{``Inbox up to date · N unread''} if unread
          messages exist, or \textit{``All caught up — no new messages''} otherwise.
          This is a specific, factual outcome statement, not a generic ``Done.'' The user
          knows the exact result without looking anywhere else.
  \end{enumerate}
  A 3-second cooldown (N5 Error Prevention) prevents accidental double-fires that
  would flood the interface with overlapping skeleton states. The button uses
  \texttt{aria-label} that dynamically updates to ``Checking for new mail\ldots''
  during the reload — supporting screen-reader users (U1 Equitable Use).
\end{lucidSolutionBox}

\textbf{Principle citations:} S3 (Informative Feedback — spin icon + ``Checking'' toast),
S4 (Closure — outcome toast after reload completes),
G1 (Closure/Reification — skeleton rows anticipate real content),
N1 (Visibility of System Status — every phase is visible),
N5 (Error Prevention — 3-second cooldown),
D6 (Affordance — circular-arrow universally signals reload),
L3 (Fitts' Law — button adjacent to label in header's F-pattern zone),
U1 (Equitable Use — dynamic aria-label for screen readers).

\begin{figure}[H]
  \centering
  \begin{minipage}[t]{0.48\textwidth}
    \centering
    \includegraphics[width=\linewidth]{reload_func_1.png}
    \caption*{\textit{Phase 1\,\&\,2 --- Spin icon + skeleton rows:}\\
    ``Checking for new mail\ldots'' toast fires immediately;\\ email list replaced by animated skeleton rows (G1, S3).}
  \end{minipage}
  \hfill
  \begin{minipage}[t]{0.48\textwidth}
    \centering
    \includegraphics[width=\linewidth]{reload_func_2.png}
    \caption*{\textit{Phase 3 --- Closure toast:}\\
    ``Inbox up to date \textperiodcentered\ 17 unread'' \\ confirms the exact outcome (S4, N1).}
  \end{minipage}
  \caption{Issue 13 --- Lucid's three-phase reload: skeleton loading state (left) and closure toast on completion (right). Gmail provides neither.}
\end{figure}


%─────────────────────────────────────────────────
\section{Issue 14 — Gmail's Stats Header Is Purely Decorative}

\begin{gmailIssueBox}{Gmail Issue: Non-Interactive Status Information}
  Gmail provides no at-a-glance overview of account state (unread count, today's
  sent count, draft count) without manual folder navigation. Even if such counts
  were displayed, clicking them in Gmail performs no action — information is
  presented but not made actionable. This creates unnecessary cognitive load:
  users who want to jump to unread mail must navigate to Inbox and scroll,
  and those who want to check today's sent mail must click the Sent folder and
  mentally filter by date.
\end{gmailIssueBox}

\begin{lucidSolutionBox}{Lucid Solution: Actionable Stat Badges with Instant Drill-Down}
  Lucid's Welcome Header displays three interactive stat badges, each triggering
  a context-appropriate action:

  \begin{itemize}[leftmargin=*, itemsep=4pt]
    \item \textbf{Unread badge} — Shows the current unread count in brand blue.
          Clicking it applies an \textit{unread-only filter} to the inbox without
          navigating away from it. The badge turns blue with a border (N1 Visibility)
          and the label changes to ``Unread \texttimes'' to signal the active state.
          A dismissable chip (``Unread only ×'') appears in the list header as an
          always-visible escape hatch (S7 Locus of Control, N3 User Control).
          Clicking the badge again, clicking the chip, or navigating away all
          revert the filter — matching users' expected reversibility (S5).

    \item \textbf{Sent badge} — Shows only \textit{today's} sent count, not a
          cumulative all-time figure. This matches the user's mental model
          (N2 Match Real World — ``How many did I send today?'' not ``How many
          have I ever sent?''). Clicking navigates to the Sent folder pre-filtered
          to today, with the list header reading ``Sent · Today (today only)''.
          A ``Today only ×'' chip clears the filter to reveal all sent mail.

    \item \textbf{Drafts badge} — Shows the current draft count. Clicking navigates
          directly to the Drafts folder — one click from the inbox instead of two
          (sidebar click).
  \end{itemize}
\end{lucidSolutionBox}

\textbf{Principle citations:} S7 (Locus of Control — user activates and reverts
filters at will), N1 (Visibility — active filter state reflected in badge colour
and chip), N3 (User Control and Freedom — escape hatch always one click away),
N2 (Match Real World — Sent badge shows today's count, not lifetime total),
N7 (Flexibility and Efficiency — one-click drill-down vs. two-click sidebar navigation),
S8 (Reduce STM Load — no need to remember folder names or navigate to find counts),
L7 (Zipf's Least Effort — maximum information with minimum clicks).

\begin{figure}[H]
  \centering
  \begin{minipage}[t]{0.32\textwidth}
    \centering
    \includegraphics[width=\linewidth]{stat_1.png}
    \caption*{\textbf{Unread badge (active)}\\
    Badge highlights blue; list header reads ``Inbox --- Unread'';
    dismissable ``Unread only \texttimes'' chip is the escape hatch (S7, N3, N1).}
  \end{minipage}
  \hfill
  \begin{minipage}[t]{0.32\textwidth}
    \centering
    \includegraphics[width=\linewidth]{stat_2.png}
    \caption*{\textbf{Sent badge --- Today filter}\\
    Navigates to Sent pre-filtered to today; header reads ``Sent \textperiodcentered\ Today (today only)''; ``Today only \texttimes'' chip clears filter (N2, L7).}
  \end{minipage}
  \hfill
  \begin{minipage}[t]{0.32\textwidth}
    \centering
    \includegraphics[width=\linewidth]{stat_3.png}
    \caption*{\textbf{Drafts badge}\\
    One click from inbox to Drafts folder; count surfaced without folder navigation (S8, N7).}
  \end{minipage}
  \caption{Issue 14 --- Actionable stat badges: unread-only filter with revert chip (left), sent-today scoped view (centre), and direct Drafts navigation (right). All three replace two-click sidebar navigation with a single badge press.}
\end{figure}


% ══════════════════════════════════════════════════════════════
\chapter{User Support Design}
% ══════════════════════════════════════════════════════════════

\section{Context-Sensitive Help (Tooltips)}

Every icon-only button in Lucid (sidebar compact mode, Top Bar icon buttons, Email Detail
action sidebar, Compose toolbar) has a \texttt{Tooltip} component that appears on hover
without disrupting the workflow. The tooltip content is written as an action description,
not a feature name: ``Compose new email'' rather than ``Compose''; ``Archive (removes from
inbox)'' rather than ``Archive''; ``Delete — moves to Trash (30 days to recover)'' rather
than ``Delete''. This satisfies User Support requirements \textbf{SP1 (Availability)},
\textbf{SP2 (Accuracy)}, and \textbf{SP3 (Unobtrusiveness)}.

\begin{figure}[H]
  \centering
  \includegraphics[width=\linewidth]{shortcut_icon.jpg}
  \caption{Context-Sensitive Help: Help icon (?) in Top Bar triggers shortcut overlay; icon buttons show tooltip on hover}
\end{figure}

\section{Keyboard Shortcut Overlay}

The \texttt{?} key and the Help icon (Top Bar) open a full-screen overlay categorising all
shortcuts into three groups. This is a \textit{quick reference} style of user support —
designed for users who are familiar with the system but need a reminder of a specific
command. The overlay dismisses with \texttt{Escape}, following the convention that any
modal overlay is escapable by the most universal exit key.

\begin{figure}[H]
  \centering
  \includegraphics[width=\linewidth]{docs.jpg}
  \caption{Keyboard Shortcut Overlay: three categories (Actions, Navigation, Search \& Help) with verb-form descriptions}
\end{figure}

\section{Settings as Transparent Control Panel}

The Settings page is both a configuration interface and a form of user support — it makes
the system's configurable capabilities visible and reachable. The ``Keyboard shortcuts''
section with descriptions, default keys, and remap badges is simultaneously a quick
reference and an interactive tool. This satisfies \textbf{SP4 (Flexibility)} — expert
users customise; casual users read.

\begin{figure}[H]
  \centering
  \includegraphics[width=\linewidth]{set1.jpg}
  \caption{Settings as Control Panel: all configurable capabilities visible and reachable in one place}
\end{figure}

\section{Compose Window Guidance}

The compose body's placeholder text doubles as inline documentation: ``Write your message…
(Ctrl+B bold, Ctrl+I italic, Ctrl+U underline)'' teaches three keyboard shortcuts in the
context where they are immediately usable — the definition of context-sensitive help.

The subject-empty warning bar with ``Send anyway'' is a \textit{wizard-style} correction
path: the system catches a potential error, explains it in one sentence, and offers the
next step in the same UI location — not in a separate help document.

\begin{figure}[H]
  \centering
  \includegraphics[width=\linewidth]{errorhand.jpg}
  \caption{Compose Window Guidance: subject-empty warning bar with ``Send anyway'' wizard-style correction path}
\end{figure}

\section{Empty State Guidance}

Every empty state in Lucid includes:
\begin{enumerate}[leftmargin=*, itemsep=1pt]
  \item An illustration (visual channel — Dual Coding Theory, P4).
  \item A plain-English title explaining the current state.
  \item A descriptive sentence explaining what goes here.
  \item An action button (where applicable) for the most relevant next step.
\end{enumerate}

This structure provides \textit{task-specific help} at the moment of confusion — without
requiring the user to navigate to a help centre.

\begin{figure}[H]
  \centering
  \includegraphics[width=\linewidth]{draft.jpg}
  \caption{Empty State Guidance: folder-specific illustration, plain-English title, description, and relevant action CTA}
\end{figure}

% ══════════════════════════════════════════════════════════════
\chapter{Accessibility and Inclusive Design}
% ══════════════════════════════════════════════════════════════

\section{Semantic HTML and ARIA}

Lucid uses semantic HTML roles throughout:
\begin{itemize}[leftmargin=*, itemsep=2pt]
  \item \texttt{<header role="banner">} — Top Bar
  \item \texttt{<aside role="navigation" aria-label="Mail navigation">} — Sidebar
  \item \texttt{<main role="main" aria-label="Email content">} — Content area
  \item \texttt{<nav role="list" aria-label="Mail folders">} — Sidebar nav
  \item \texttt{<div role="list" aria-label="Email list">} — Email list
  \item \texttt{<div role="dialog" aria-label="Compose email">} — Compose window
  \item \texttt{<button role="switch" aria-checked>} — Toggle switches in Settings
\end{itemize}

\begin{figure}[H]
  \centering
  \includegraphics[width=\linewidth]{inbox.jpg}
  \caption{Semantic HTML and ARIA: sidebar with \texttt{role="navigation"}, email list with \texttt{role="list"}, active item \texttt{aria-current="page"}}
\end{figure}

\section{Colour Independence}

No information is conveyed by colour alone. All read/unread distinctions are conveyed by
both colour (\texttt{text-primary} vs. \texttt{text-secondary}) and font weight (bold vs.
normal). Danger actions are both red and labelled with explicit action text.
Custom tag colours in compact sidebar mode are supplemented by the first letter of the
tag name — readable without colour.

\begin{figure}[H]
  \centering
  \includegraphics[width=\linewidth]{preview1.jpg}
  \caption{Colour Independence: read/unread conveyed by bold weight \emph{and} contrast --- no colour-only channel}
\end{figure}

\section{Keyboard Navigability}

All interactive elements are reachable by Tab. The keyboard shortcut system provides
an alternative navigation path to every major function without requiring mouse use.
Form fields have logical tab order within each form group (To → Cc → Bcc → Subject →
Body → Send in the Compose Window).

\begin{figure}[H]
  \centering
  \includegraphics[width=\linewidth]{shortcut.jpg}
  \caption{Keyboard Navigability: full shortcut overlay providing keyboard paths to every major function}
\end{figure}

\section{Relative Sizing}

The entire application uses \texttt{rem} units derived from a root font size set by the
zoom widget (\texttt{document.documentElement.style.fontSize}). When the user increases
zoom from 100\% to 150\%, all text, padding, margins, and icon sizes scale proportionally.
No hardcoded \texttt{px} measurements exist in layout-critical code paths.

\begin{figure}[H]
  \centering
  \includegraphics[width=\linewidth]{resize.jpg}
  \caption{Relative Sizing: zoom widget scales all text, padding, and icon sizes proportionally via root \texttt{rem} unit}
\end{figure}

\section{Screen Reader Support}

\texttt{aria-label} attributes are provided for:
\begin{itemize}[leftmargin=*, itemsep=1pt]
  \item Every icon button (``Go to inbox'', ``Compose new email'', ``Archive selected'')
  \item Every navigation item (\texttt{aria-current="page"} for active items)
  \item Every toggle switch (\texttt{aria-checked={value}})
  \item Every form input (``Email subject'', ``Reply body'', ``Tag name'')
  \item Every avatar (sender name as aria-label)
  \item Decorative overlay (\texttt{aria-hidden="true"} on the compose backdrop)
\end{itemize}

\begin{figure}[H]
  \centering
  \includegraphics[width=\linewidth]{expand_closure.jpg}
  \caption{Screen Reader Support: compact and expanded sidebar modes both expose \texttt{aria-label} on every navigation item}
\end{figure}

% ══════════════════════════════════════════════════════════════
\chapter{Conclusion}
% ══════════════════════════════════════════════════════════════

\section{Design Principles Summary}

The following table summarises the mapping of every major Lucid feature to its primary
governing principle:

\begin{table}[H]
\centering
\small
\begin{tabularx}{\textwidth}{lXl}
  \toprule
  \textbf{Feature} & \textbf{Design Decision} & \textbf{Primary Principle}\\
  \midrule
  Night sky landing page       & Visceral positive first impression                  & P3, P1\\
  LUCID wordmark colours       & Google colour transfer — trust and familiarity       & P6, L8\\
  Resizable sidebar            & User controls layout, drag-to-resize                & S7, D4, D6\\
  Search always visible        & Feature always in primary visual field               & D1, W4\\
  Recent search history        & System remembers so user doesn't                    & S8, N6\\
  Welcome header with stats    & At-a-glance state without navigation                & S8, N1\\
  Skeleton loading rows        & Closure replaces blank loading state                & G1, P3\\
  Context-specific empty states & Folder-aware guidance and CTA                      & N2, N9\\
  Bulk action bar              & Context-aware actions, clear selection mode exit    & S7, N3, G2\\
  Email detail two-column      & Actions proximate to content; Figure-Ground focus   & G3, G7\\
  Trash/Archive context banner & Closure: user always knows email's state            & S4, S8\\
  Smart replies                & Least effort for common responses                   & L7, N7\\
  Autosave every 30s           & Never lose draft work                               & L10(1), S3\\
  Subject-empty guard          & Prevent most common compose error                   & S6, N5\\
  Send button disabled (no To) & Constrain to valid actions only                     & D3, N5\\
  Toast with Undo              & Reversibility for all destructive actions            & S5, N3\\
  Keyboard Z to undo           & Shortcut for expert reversal                        & S2, S5\\
  Keyboard shortcuts system    & Expert navigation without mouse                     & S2, N7, L3\\
  Remappable shortcuts         & Respect prior practice curve learning               & L5, S7\\
  Zoom widget in Top Bar       & Per-app scaling independent of browser              & U7, U2\\
  Dark mode theme              & Reduce eye strain; user preference                  & U2, S7\\
  Three density modes          & Adapts to scanning vs. reading workflows            & N7, U2\\
  Colour-coded tag pills       & Dual-channel identification (colour + letter)       & G6, U4\\
  Password shown by default    & Prevent typos without security theatre               & N5, W6\\
  Guest demo mode              & Zero-friction evaluation path                       & W2, P5\\
  Social proof + testimonials  & Build trust before commitment                       & P5, P2\\
  Footer: Privacy Terms Help   & Standard web conventions                            & L8, W9\\
  Reload button (3-phase)      & Spin icon + skeleton rows + closure toast          & S3, S4, G1, N1\\
  Stat badge — Unread filter   & One-click inbox drill-down; dismissable chip to revert & S7, N3, N1\\
  Stat badge — Sent today      & Badge count = today only; filter pre-applied on click   & N2, L7, S8\\
  Stat badge — Drafts nav      & One-click jump to Drafts; counts surface without nav & S8, N7\\
  \bottomrule
\end{tabularx}
\caption{Master feature-to-principle mapping}
\end{table}

\section{Final Reflection}

Lucid Mail demonstrates that the gap between a feature-rich product like Gmail and a
\textit{usable} product is not one of capability, but of design intent. Gmail has every
feature Lucid has, and hundreds more. Yet by adding features without principled constraints,
Gmail has accumulated cognitive debt: a Promotions tab that buries real emails, a fixed
sidebar that cannot accommodate power users, a compose window whose size is non-negotiable,
and no at-a-glance dashboard for the user's own account state.

Lucid's approach — one principle, one feature, every time — produces an interface that
is objectively narrower in feature count but subjectively easier to use. Tesler's Law
reminds us that the complexity did not disappear; it was relocated from the user's working
memory to the developer's implementation. That is the correct trade-off for a well-designed
human-centred system.

The 40 years of HCI research encoded in Shneiderman, Nielsen, Norman, Gestalt, Raskin,
Krug, and their contemporaries are not academic exercises. They are empirical findings
about how human cognition works under load. Lucid is an argument in code that
following them systematically produces a product that users describe not as ``simple''
or ``limited'' but as ``clean'', ``fast'', and — most importantly — ``it just works.''

\begin{figure}[H]
  \centering
  \includegraphics[width=\linewidth]{mainpage.jpg}
  \caption{Lucid Mail landing page --- the starting surface that sets the first impression}
\end{figure}

\begin{figure}[H]
  \centering
  \includegraphics[width=\linewidth]{starred.jpg}
  \caption{Lucid Mail application --- inbox to email detail view, every decision anchored to a named HCI principle}
\end{figure}

% ══════════════════════════════════════════════════════════════
\chapter{Loading States \& Onboarding}
\addcontentsline{toc}{chapter}{Loading States \& Onboarding}
% ══════════════════════════════════════════════════════════════

Loading and first-run experiences are among the most overlooked surfaces in email client
design, yet they set the user's entire mental model of a product. Lucid implements three
distinct loading treatments, each grounded in a specific HCI principle.

% ── 12.1 ───────────────────────────────────────────────────────
\section{Login Loading Screen — Animated Wordmark}

\begin{tcolorbox}[enhanced, breakable, colback=brandlight, colframe=principleblue, left=6pt, right=6pt, top=5pt, bottom=5pt, boxrule=1.2pt, arc=4pt]
  \textbf{Principle anchor:} Shneiderman S3 (Informative Feedback) · Gestalt G1 (Closure)
  · Aesthetic Usability Effect P1
\end{tcolorbox}

Immediately after the user authenticates, Lucid displays a full-screen interstitial
featuring the ``LUCID'' wordmark rendered in a neutral grey. A left-to-right colour sweep
progressively fills each letter with its signature Google colour
(\textcolor{brand}{\#4285F4} · \textcolor{danger}{\#EA4335} · \textcolor{warning}{\#FBBC04} · \textcolor{success}{\#34A853})
as a rainbow progress bar beneath it advances in parallel.

\textbf{Design rationale.}
\begin{itemize}
  \item \textit{Informative Feedback (S3):} The advancing colour fill gives the user
        continuous, proportional feedback about progress — not a spinner that communicates
        nothing about how long remains.
  \item \textit{Closure (G1):} Partial colouring creates a perceptual ``gap'' the brain
        actively wants to close. Users perceive the incomplete state as temporary rather than
        broken, reducing impatience.
  \item \textit{Aesthetic Usability Effect (P1):} The on-brand animation leaves a premium
        first impression that colours the user's evaluation of every subsequent interaction
        (Aesthetic Usability Effect, Kurosu \& Kashimura, 1995).
  \item \textit{Familiar pattern — novel execution:} Progress bars are a known affordance
        (Jakob's Law, L8). Threading them through the product wordmark is distinctive without
        being disorienting.
\end{itemize}

The screen fades out over 480 ms once the animation completes, transitioning smoothly to
the skeleton or inbox without a jarring cut.

\begin{figure}[H]
  \centering
  
\includegraphics[width=\linewidth]{loading.jpg}
  \caption{Login Loading Screen: LUCID wordmark colour sweep with rainbow progress bar beneath}
\end{figure}

% ── 12.2 ───────────────────────────────────────────────────────
\section{Inbox Skeleton — Full-Layout Shimmer}

\begin{principleBox}{Principle Anchor}
  \textbf{Principle anchor:} Gestalt G1 (Closure) · Shneiderman S3 (Feedback) ·
  Norman N1 (Visibility of System Status)
\end{principleBox}

When the main application mounts after login, Lucid renders a pixel-accurate skeleton of
the complete three-column layout — topbar, sidebar navigation, and a column of twelve email
rows — for 1.3 seconds before the real data appears. Every skeleton element pulses at
1.6 s using a \texttt{sk-pulse} CSS keyframe animation that oscillates opacity between 1
and 0.45.

\textbf{Why a full-layout skeleton outperforms a spinner:}
\begin{enumerate}
  \item \textbf{Spatial priming.} The skeleton occupies exactly the space the real content
        will fill. Users' eyes land in the right place before data arrives, eliminating the
        jarring reflow that occurs when a spinner is replaced by a full layout (Luke
        Wroblewski, ``Mobile First'', 2011).
  \item \textbf{Perceived performance.} Skeleton screens consistently score 20–30\% faster
        on perceived-load tests than equivalent spinner-based loading (Facebook, 2013;
        Lincoln Loop, 2015). Lucid's 1.3 s real-data delay is masked almost entirely by the
        skeleton's active appearance.
  \item \textbf{Closure (G1).} Partially rendered shapes trigger the Gestalt law of closure:
        the brain interprets incomplete rectangles as email rows and completes the picture,
        mentally ``pre-loading'' the inbox before it arrives.
  \item \textbf{Structural hints.} The skeleton sidebar reproduces the item count and active
        highlight of the real sidebar, communicating the app's information architecture
        before a single byte of email data has loaded.
\end{enumerate}

The skeleton uses the shared \texttt{.skeleton} CSS class (background \texttt{var(--gray-100)},
border-radius \texttt{var(--radius-sm)}) for visual consistency with the in-list loading
rows already present in the email list component.

\begin{figure}[H]
  \centering
  \includegraphics[width=\linewidth]{inbox.jpg}
  \caption{Inbox Skeleton: pixel-accurate three-column layout (topbar, sidebar, email rows) pulsing before real data arrives}
\end{figure}

% ── 12.3 ───────────────────────────────────────────────────────
\section{Onboarding Tutorial — Blurred Backdrop Popup}

\begin{principleBox}{Principle Anchor}
  \textbf{Principle anchor:} Jakob's Law L8 · Shneiderman S8 (Reduce STM Load) ·
  Cialdini P2 (Liking) · Norman N9 (Help \& Documentation)
\end{principleBox}

First-time users — identified by the absence of a \texttt{lucid\_onboarded\_v1}
\texttt{localStorage} key — are presented with a five-slide modal tutorial immediately after
the skeleton resolves. The tutorial appears over a 10 px Gaussian-blur backdrop
(\texttt{backdrop-filter: blur(10px)}), keeping the inbox visible but defocused so the
user retains spatial context while reading the tutorial.

\subsection*{Slide structure}

\begin{center}
\begin{tabular}{@{}lll@{}}
  \toprule
  \textbf{Slide} & \textbf{Title} & \textbf{Core message} \\
  \midrule
  1 & Welcome to Lucid        & Side-by-side comparison: what's the same vs.\ what's better \\
  2 & Familiar layout         & Three-column layout identical to Gmail — zero learning curve \\
  3 & Keyboard shortcuts      & C / E / Z / ? — all the power-user keys at a glance \\
  4 & Undo anything, anytime  & Contrast with Gmail's 30-second undo-send limitation \\
  5 & You're all set          & Confirmation + hint to press ? for the full shortcut overlay \\
  \bottomrule
\end{tabular}
\end{center}

Each slide pairs a custom SVG illustration with concise text. Inline SVG was chosen over
raster images to guarantee pixel-perfect rendering at all display densities without
network round trips.

\textbf{Design rationale.}
\begin{itemize}
  \item \textit{Jakob's Law (L8):} Framing Lucid relative to Gmail explicitly transfers the
        user's existing mental model. Slide 1 lists every familiar affordance alongside the
        specific improvement Lucid makes — reducing re-learning to zero.
  \item \textit{Reduce STM Load (S8):} Five short slides present one idea each. Chunking
        (Miller's Law) keeps working memory within its 7 ± 2 capacity.
  \item \textit{Blurred backdrop — figure/ground (G3/G7):} The defocused inbox behind the
        modal establishes the modal as figure and the inbox as ground, creating clear visual
        hierarchy without hiding the spatial context of where the user has arrived.
  \item \textit{Liking (P2):} The illustrated, conversational tone of the tutorial increases
        affinity with the product before the user has written a single email (Cialdini's
        Liking principle — people respond positively to things that feel personal and warm).
  \item \textit{Skip affordance:} A ``Skip'' button is always visible in the top-right corner
        of the modal, honouring user autonomy (N3 User Control) and avoiding the dark pattern
        of forced onboarding flows.
  \item \textit{Progressive dots + Back navigation:} Dot indicators (current slide
        highlighted in brand blue, others in \texttt{var(--gray-200)}) communicate position
        within the sequence without requiring the user to count slides. Back navigation
        supports exploration rather than forcing a linear path.
  \item \textit{One-time display:} The flag is stored in \texttt{localStorage} rather than
        a session variable, so returning users are never shown the tutorial again — respecting
        their expertise and eliminating the patronising repetition that plagues many web apps.
\end{itemize}

\begin{figure}[H]
  \centering
  \includegraphics[width=\linewidth]{login_gmail.jpg}
  \caption{Onboarding Slide 1: Gmail vs.\ Lucid comparison --- familiar affordances listed alongside specific improvements}
\end{figure}

\begin{figure}[H]
  \centering
  \includegraphics[width=\linewidth]{feedback4.jpg}
  \caption{Onboarding Slide 4: Undo anything, anytime --- contrasting Lucid's persistent undo with Gmail's 30-second window}
\end{figure}

% ── 12.4 ───────────────────────────────────────────────────────
\section{Loading State Sequence Summary}

The three loading surfaces form a deliberate \textbf{temporal funnel}:

\begin{center}
\begin{tabular}{@{}clll@{}}
  \toprule
  \textbf{Phase} & \textbf{Surface} & \textbf{Duration} & \textbf{User state} \\
  \midrule
  1 & Login loading screen (wordmark fill) & $\approx$2.2 s & Authenticating     \\
  2 & Inbox skeleton (full layout shimmer) & $\approx$1.3 s & App bootstrapping  \\
  3 & Onboarding tutorial (blurred modal)  & User-controlled & First-run only     \\
  \bottomrule
\end{tabular}
\end{center}

Each phase has a distinct visual language — brand wordmark animation, structural skeleton,
contextual modal — so the user always knows \textit{why} they are waiting and \textit{what}
comes next. Together they transform an unavoidable delay into a designed moment that builds
brand familiarity, transfers mental models, and sets expectations before the user sends
their first email.

% ══════════════════════════════════════════════════════════════
\chapter*{References}
\addcontentsline{toc}{chapter}{References}
% ══════════════════════════════════════════════════════════════

\begin{itemize}[leftmargin=*, itemsep=6pt]
  \item Shneiderman, B. (1987). \textit{Designing the User Interface}. Addison-Wesley.
  \item Nielsen, J. (1994). ``10 Usability Heuristics for User Interface Design.''
        Nielsen Norman Group.
  \item Norman, D. A. (1988). \textit{The Design of Everyday Things}. Basic Books.
  \item Miller, G. A. (1956). ``The Magical Number Seven, Plus or Minus Two.''
        \textit{Psychological Review}, 63(2), 81--97.
  \item Fitts, P. M. (1954). ``The Information Capacity of the Human Motor System.''
        \textit{Journal of Experimental Psychology}, 47(6), 381--391.
  \item Hick, W. E. (1952). ``On the Rate of Gain of Information.''
        \textit{Quarterly Journal of Experimental Psychology}, 4(1), 11--26.
  \item Gestalt Psychology — Wertheimer, M.; Koffka, K.; Köhler, W. (1920s–1930s).
        Various publications in \textit{Psychologische Forschung}.
  \item Krug, S. (2000). \textit{Don't Make Me Think}. New Riders.
  \item Raskin, J. (2000). \textit{The Humane Interface}. Addison-Wesley.
  \item Cialdini, R. B. (1984). \textit{Influence: The Psychology of Persuasion}. Harper Business.
  \item Tesler, L. (1980). ``A Personal Computer for Children of All Ages.''
        Proceedings of ACM 1980.
  \item Pareto, V. (1896). \textit{Cours d'Économie Politique}. Droz.
  \item Nielsen, J. (2014). ``Password Masking: Privacy, Not Security.''
        Nielsen Norman Group.
  \item Apple Inc. (2023). \textit{Human Interface Guidelines}. Apple Developer.
  \item Story, M. F. (1998). ``Maximising Usability: The Principles of Universal Design.''
        \textit{Assistive Technology}, 10(1), 4--12.
  \item HCI Exam Reference Guide (2026). Course slide compilation.
\end{itemize}

\section{Useful Links}

\begin{itemize}[leftmargin=*, itemsep=6pt]
  \item \textbf{Lucid Mail --- GitHub Repository}\\[2pt]
        \href{https://github.com/sxdxde/Lucid}{\texttt{https://github.com/sxdxde/Lucid}}
\end{itemize}

  % ── End Course & Team Details ──────────────────────────────
  \clearpage
  \thispagestyle{empty}
\pagecolor{dark}
\color{white}
\vspace*{\fill}
\begin{center}
  \begin{tcolorbox}[
    enhanced, sharp corners,
    colback=white!5!dark, colframe=white!15!dark,
    width=0.65\textwidth,
    left=28pt, right=28pt, top=28pt, bottom=28pt,
  ]
    \centering\color{white!80!dark}
    {\large\bfseries\color{white} Course Details}\\[14pt]
    {\normalsize\color{white!85!dark} Human Computer Interaction (HCI)}\\[6pt]
    {\normalsize\color{white!55!dark} CS5015}\\[24pt]
    {\large\bfseries\color{white} Team Members}\\[14pt]
    {\normalsize\color{white!85!dark} Sudarshan S}\quad
    {\normalsize\color{white!45!dark} CS23B2007}\\[10pt]
    {\normalsize\color{white!85!dark} R Abinav}\quad
    {\normalsize\color{white!45!dark} ME23B1004}
  \end{tcolorbox}
\end{center}
\vspace*{\fill}
\clearpage

\end{document}
